\documentclass[twoside,11pt]{article}

% JMLR Style
\usepackage{jmlr2e}
\setcitestyle{numbers,sort&compress}

% Packages
\usepackage{amsmath,amssymb,amsfonts}
\usepackage{bm}
\usepackage{listings}
\usepackage{graphicx}
\usepackage{float}
\usepackage{booktabs}
\usepackage{caption}
\usepackage{subcaption}
\usepackage{algorithm}
\usepackage{algorithmic}
\usepackage{hyperref}
\usepackage{xcolor}

% Custom macros
\newcommand{\R}{\mathbb{R}}
\newcommand{\qed}{\hfill$\square$}

% Colors
\definecolor{successgreen}{RGB}{0,128,0}
\definecolor{failred}{RGB}{200,0,0}
\definecolor{warningorange}{RGB}{255,140,0}

\begin{document}

\jmlrheading{XX}{2025}{1-48}{1/25}{9/25}{Bonet Chaple}
\ShortHeadings{LLMs as Interfaces to Symbolic Discovery}{Bonet Chaple}\citep{koza1994genetic}\citep{kambhampati2024llms}
\firstpageno{1}

\title{Large Language Models as Interfaces to Symbolic Discovery: \\ A Hybrid Architecture for Reliable Analytical Expression Discovery}\citep{koza1994genetic}\citep{lewkowycz2022minerva}

\author{\name Ruperto Pedro Bonet Chaple \email ruperto.bonet@modelphysmat.com \\
       \addr Independent Researcher \\
             London, United Kingdom
}

\editor{TBD}

\maketitle

\begin{abstract}
Can Large Language Models discover novel analytical expressions, or do they primarily reproduce patterns from training data? We address this question through systematic empirical evaluation across 131 test cases spanning classical physics and decentralized finance domains. Our results reveal that pure LLM-based discovery exhibits stark domain-dependent performance: 95\% success on classical formulas versus 45\% on specialized financial mathematics, with catastrophic extrapolation errors (412-847\% versus 8-23\% for symbolic methods). We prove that autoregressive language models are fundamentally reproductive systems, with generation probability decaying exponentially for expressions distant from training data. This limitation, confirmed by recent systematic evaluations \citep{xu2025fundamental,du2025evaluating}, necessitates hybrid approaches.

Motivated by these limitations, we introduce \textsc{HypatiaX}, a hybrid symbolic-neural architecture where LLMs serve as \emph{interfaces} to symbolic discovery engines rather than discovery engines themselves. Our system combines evolutionary symbolic regression with multi-layer validation (dimensional analysis, domain constraints, execution verification) and optional LLM initialization for efficiency. Across five experimental campaigns, \textsc{HypatiaX} achieves 88.9-100\% success rates with 8-23\% extrapolation error and 100\% functional correctness on ground truth tests, compared to 20-40\% correctness for pure LLM approaches. An LLM-guided variant reduces discovery time by 82\% while preserving all symbolic guarantees. Our key contribution is demonstrating that reliable analytical discovery requires symbolic methods under physics constraints, with LLMs providing interpretation rather than mathematical reasoning\citep{koza1994genetic}\citep{la2021contemporary}\citep{kambhampati2024llms}\citep{lewkowycz2022minerva}.
\end{abstract}

\begin{keywords}
symbolic regression, large language models, scientific machine learning, hybrid AI systems, analytical expression discovery, physics-informed machine learning\citep{koza1994genetic}\citep{molnar2020interpretable}\citep{rudin2019stop}\citep{la2021contemporary}\citep{lewkowycz2022minerva}\citep{virtanen2020scipy}
\end{keywords}

\section{Introduction}

Large Language Models (LLMs) have demonstrated remarkable capabilities in scientific tasks, from literature synthesis to mathematical problem-solving \citep{brown2020language,wei2022chain,openai2023gpt4}. Recent work suggests these models can recover closed-form equations from data \citep{romera2023discovering}, raising a fundamental question: \textbf{Do LLMs perform genuine analytical discovery, or do they primarily reproduce memorized patterns?}

This distinction has profound implications for trustworthy AI in science\citep{koza1994genetic}. Genuine discovery requires:
\begin{enumerate}
\item \textbf{Generalization beyond training data}: Finding expressions for phenomena not encountered during pre-training\citep{liu2020understanding}\citep{neyshabur2017exploring}\citep{zhang2021understanding}
\item \textbf{Physical consistency}: Respecting dimensional analysis, conservation laws, and domain constraints
\item \textbf{Extrapolation capability}: Producing predictions that remain valid outside the training distribution
\item \textbf{Failure transparency}: Detecting and rejecting incorrect expressions rather than producing plausible but wrong results
\end{enumerate}

\subsection{Motivating Evidence}

Consider two formula discovery tasks:\citep{koza1994genetic}

\paragraph{Classical Physics (Ohm's Law):} Given voltage-current-resistance data, discover $V = IR$.

\paragraph{DeFi Risk (Value-at-Risk):} Given return distribution data, discover $\text{VaR}_\alpha = -\mu + \sigma \Phi^{-1}(\alpha)$.

When we evaluated state-of-the-art LLMs \citep{openai2023gpt4,anthropic2024claude} on 40 such problems:
\begin{itemize}
\item \textbf{Classical formulas}: 95\% success rate, perfect $R^2 \approx 1.0$
\item \textbf{Financial formulas}: 45\% success rate, frequent failures with $R^2 < 0$
\item \textbf{Extrapolation}: 847\% error for LLMs versus 23\% for symbolic methods\citep{kambhampati2024llms}
\item \textbf{Functional correctness}: 20-40\% for LLMs versus 88\citep{kambhampati2024llms}.9-100\% for our hybrid approach
\end{itemize}

This bimodal behavior suggests LLMs excel at reproducing well-known relationships but struggle with specialized domains requiring distributional reasoning or novel functional forms\citep{kambhampati2024llms}\citep{lewkowycz2022minerva}.

\subsection{Contributions}

This paper makes four primary contributions:

\begin{enumerate}
\item \textbf{Theoretical characterization} of LLMs as reproductive systems (Theorem~\ref{thm:llm_reproductive}), with proof that autoregressive generation probability decays exponentially for expressions distant from training data\citep{kambhampati2024llms}.

\item \textbf{Comprehensive empirical evaluation} demonstrating domain-dependent LLM performance across 131 tests: 95\% success on classical physics versus 45\% on decentralized finance, with detailed failure mode taxonomy\citep{kambhampati2024llms}.

\item \textbf{Hybrid architecture} (\textsc{HypatiaX}) combining evolutionary symbolic regression with multi-layer validation and optional LLM initialization, achieving 88.9-100\% success rates with 8-23\% extrapolation error\citep{la2021contemporary}\citep{kambhampati2024llms}.

\item \textbf{Design principles} for reliable analytical discovery: symbolic methods for mathematical reasoning, LLMs for natural language interpretation and warm-starting search\citep{koza1994genetic}\citep{lewkowycz2022minerva}\citep{kambhampati2024llms}.
\end{enumerate}

\subsection{Paper Organization}

We begin with related work (Section~\ref{sec:related}), then present empirical evidence of LLM limitations (Section~\ref{sec:llm_limitations}) to motivate our approach. Section~\ref{sec:theory} provides theoretical analysis, followed by the \textsc{HypatiaX} architecture (Section~\ref{sec:architecture}). Experimental validation (Section~\ref{sec:experiments}) demonstrates performance across five campaigns, and we conclude with discussion and limitations (Section~\ref{sec:discussion})\citep{kambhampati2024llms}.

\section{Related Work}
\label{sec:related}

\subsection{Symbolic Regression}\citep{la2021contemporary}

Symbolic regression discovers analytical expressions through combinatorial search over mathematical operators. Classical genetic programming approaches \citep{koza1992genetic,schmidt2009distilling} evolve expression trees using fitness-based selection. Modern frameworks like PySR \citep{cranmer2023interpretable} incorporate Pareto optimization to balance accuracy and complexity, physics-informed operators, and dimensional constraints. Recent neural-guided approaches \citep{petersen2021deep,udrescu2020aifeynman} combine deep learning with symbolic search, while \citet{cranmer2023symbolic} provides a comprehensive review of symbolic regression in the modern era. Cutting-edge methods integrate LLMs with symbolic regression \citep{taskin2026knowledge}, develop neural-symbolic models for physics \citep{ying2025physe2e}, and apply physics-guided deep SR to engineering problems \citep{zhang2026formula,senetakis2026symbolic}.

Our work builds on these foundations by integrating symbolic regression with LLM initialization and multi-layer validation, achieving superior performance on specialized domains while maintaining formal guarantees\citep{la2021contemporary}\citep{kambhampati2024llms}.

\subsection{LLMs for Scientific Discovery}\citep{koza1994genetic}\citep{virtanen2020scipy}\citep{kambhampati2024llms}

Recent work has explored LLMs for hypothesis generation \citep{wang2023scientific}, literature mining \citep{liu2023summary}, and mathematical reasoning \citep{wei2022chain,zhang2025survey}. Advanced reasoning techniques include adversarial reinforcement learning \citep{liu2025adversarial}, memory-augmented agents \citep{yan2026memory}, and dynamic task adaptation \citep{wang2025reasoning}. The transformer architecture \citep{vaswani2017attention} underpins these capabilities, though recent studies \citep{mirzadeh2024reasoning,dziri2023faith,xu2025fundamental,tian2025reasoning} reveal fundamental limitations in compositional reasoning and extrapolation. Systematic evaluations \citep{du2025evaluating,chen2025epistemic} demonstrate consistent performance gaps between LLM capabilities and genuine scientific discovery. Romera et al. \citep{romera2023discovering} demonstrate that LLMs can recover known physical laws from textual descriptions and data patterns.

\textbf{Key difference}: Romera et al. evaluate on well-established physics formulas likely present in training data. We systematically test on specialized domains (decentralized finance, advanced statistics) requiring genuine extrapolation, revealing fundamental limitations in pure LLM approaches\citep{kambhampati2024llms}.

\subsection{Hybrid AI Systems}

Neurosymbolic AI combines neural networks with symbolic reasoning \citep{garcez2019neural,kautz2020third}. In scientific computing, physics-informed neural networks \citep{raissi2019physics,karniadakis2021physics} embed differential equations as constraints, while neural operators \citep{lu2021learning,li2020fourier} learn mappings between function spaces. Recent surveys \citep{li2026physics} demonstrate the expanding scope of physics-informed AI across complex systems. \citet{udrescu2020ai} uses neural networks to guide symbolic search.

Our architecture differs by using LLMs purely as interfaces (initialization + interpretation) rather than as components of the discovery process, maintaining symbolic guarantees while leveraging LLM strengths\citep{koza1994genetic}\citep{kambhampati2024llms}.

\section{Empirical Evidence of LLM Limitations}\citep{kambhampati2024llms}
\label{sec:llm_limitations}

To understand whether LLMs can perform genuine analytical discovery, we conducted systematic evaluation of pure LLM-based formula generation without symbolic constraints, analytical solvers, or post-generation correction\citep{koza1994genetic}\citep{kambhampati2024llms}.

\subsection{Experimental Design}

\paragraph{Test Domains:} We evaluated across two complementary sets:

\begin{itemize}
\item \textbf{Classical Scientific (20 tests)}: Materials Science (4), Fluid Dynamics (4), Thermodynamics (4), Mechanics (4), Chemistry (4)\citep{virtanen2020scipy}
\item \textbf{Decentralized Finance (20 tests)}: Automated Market Makers (4), Value-at-Risk (4), Liquidity (4), Expected Shortfall (4), Liquidation (4)
\end{itemize}

\paragraph{Protocol:} For each test, we provided Claude Sonnet 4 with:
\begin{itemize}
\item Domain context and variable descriptions with physical units
\item Sample data statistics (mean, standard deviation, range, correlations)
\item Request for closed-form analytical expression and executable Python code
\item No symbolic regression, constraint enforcement, or correction\citep{la2021contemporary}
\end{itemize}

\paragraph{Evaluation Metrics:}
\begin{itemize}
\item $R^2$ coefficient of determination on held-out test data
\item Extrapolation error: performance on data 2$\times$ outside training range
\item Functional correctness: whether discovered expression matches known ground truth
\item Failure mode classification: execution errors, dimensional inconsistencies, statistical misuse
\end{itemize}

\subsection{Baseline Results: Bimodal Performance}

Table~\ref{tab:baseline_llm} reveals striking domain-dependent behavior\citep{kambhampati2024llms}.

\begin{table}[htbp]
\centering
\caption{Pure LLM Performance Across Domains (40 tests, Claude Sonnet 4)}
\label{tab:baseline_llm}
\begin{tabular}{@{}lcccc@{}}
\toprule
\textbf{Domain Category} & \textbf{Tests} & \textbf{Success} & \textbf{Mean $R^2$} & \textbf{Correct Form} \\
\midrule
\multicolumn{5}{c}{\textit{Classical Scientific Domains}} \\\citep{virtanen2020scipy}
\midrule
Materials Science & 4 & 100\% & 1.000 & 4/4 \\
Fluid Dynamics & 4 & 100\% & 1.000 & 4/4 \\
Thermodynamics & 4 & 100\% & 1.000 & 4/4 \\
Mechanics & 4 & 100\% & 1.000 & 4/4 \\
Chemistry & 4 & 75\% & 0.750 & 3/4 \\
\midrule
\textbf{Classical Average} & \textbf{20} & \textbf{95\%} & \textbf{0.975} & \textbf{19/20 (95\%)} \\
\midrule
\multicolumn{5}{c}{\textit{Decentralized Finance Domains}} \\
\midrule
AMM Mechanics & 4 & 75\% & 0.817 & 3/4 \\
Risk (VaR) & 4 & 75\% & $-1.813$ & 2/4 \\
Liquidity & 4 & 50\% & --- & 2/4 \\
Expected Shortfall & 4 & 25\% & $-4.120$ & 1/4 \\
Liquidation & 4 & 0\% & --- & 0/4 \\
\midrule
\textbf{DeFi Average} & \textbf{20} & \textbf{45\%} & \textbf{---} & \textbf{8/20 (40\%)} \\
\midrule
\midrule
\textbf{Overall} & \textbf{40} & \textbf{70\%} & \textbf{0.603} & \textbf{27/40 (67.5\%)} \\
\bottomrule
\end{tabular}
\end{table}

\paragraph{Key Finding:} Pure LLMs achieve near-perfect performance on classical formulas present in training corpora but fail catastrophically on specialized domains requiring distributional reasoning or novel functional forms\citep{kambhampati2024llms}\citep{lewkowycz2022minerva}.

\subsection{Failure Mode Taxonomy}

These failure modes align with known limitations of transformer-based models in systematic generalization \citep{lake2017building} and extrapolation \citep{zhang2024extrapolation}. Recent work on multimodal reasoning \citep{malkinski2025multimodal} and argumentative reasoning \citep{gould2025argumentative} further exposes structural deficits in current architectures.

We categorize LLM failures into five distinct types:\citep{kambhampati2024llms}

\begin{table}[htbp]
\centering
\caption{Pure LLM Failure Modes (Ordered by Severity)}
\label{tab:failure_modes}
\begin{tabular}{@{}lcp{7cm}@{}}
\toprule
\textbf{Failure Type} & \textbf{Frequency} & \textbf{Description} \\
\midrule
\textbf{Silent Semantic Errors} & High & Plausible expressions with high $R^2$ in training range but incorrect functional form; produces wrong predictions when deployed \\
\midrule
\textbf{Distributional Reasoning} & High & Misuse of statistical functions (wrong tail for VaR, incorrect scipy API calls, sign convention errors) \\\citep{lewkowycz2022minerva}
\midrule
\textbf{Unit/Scale Inconsistency} & Medium & Correct functional form but missing unit conversions (percentage vs. decimal, degrees vs. radians) \\
\midrule
\textbf{Non-executable Output} & Medium & Tutorial text or pseudocode instead of executable function \\
\midrule
\textbf{Incomplete Construction} & Medium & Multi-step derivations with only partial formula or skipped calculations \\
\bottomrule
\end{tabular}
\end{table}

\subsection{Case Studies}

\subsubsection{Case Study 1: Silent Semantic Error (Impermanent Loss)}

\paragraph{Ground Truth:}
\begin{equation}
\text{IL}_{\%} = \left(\frac{2\sqrt{r}}{1+r} - 1\right) \times 100
\end{equation}

\paragraph{LLM Output:}
\begin{verbatim}
def formula(price_ratio):
    r = price_ratio
    return (2 * math.sqrt(r)) / (1 + r) - 1  # Missing ×100
\end{verbatim}

\paragraph{Analysis:}
\begin{itemize}
\item Training $R^2 = 0.98$ (appears excellent)
\item Predictions systematically wrong by 100$\times$
\item Error undetectable without dimensional analysis
\item A trader using this formula would catastrophically underestimate risk
\end{itemize}

This demonstrates why $R^2$ alone is insufficient for validation.

\subsubsection{Case Study 2: Distributional Reasoning Failure (Value-at-Risk)}\citep{lewkowycz2022minerva}

Value-at-Risk is a standard risk metric in finance \citep{rockafellar2000optimization,danielsson2017forecasting}, requiring precise understanding of tail distributions.

\paragraph{Ground Truth (95\% VaR):}
\begin{equation}
\text{VaR}_{95} = -\mu + \sigma \cdot \Phi^{-1}(0.95) = -\mu + 1.645\sigma
\end{equation}

\paragraph{LLM Output:}
\begin{verbatim}
def formula(portfolio_value, daily_volatility):
    mu = 0.0
    z_score = -1.645  # WRONG SIGN
    return -mu + z_score * daily_volatility * portfolio_value
\end{verbatim}

\paragraph{Analysis:}
\begin{itemize}
\item Used negative z-score, inverting all predictions
\item Result: $R^2 = -10.05$ (catastrophic failure)
\item Statistical convention ambiguity confused the LLM\citep{kambhampati2024llms}
\item Represents semantic knowledge failure, not computational error
\end{itemize}

VaR measures loss magnitude (positive tail), but some formulations use negative returns. The LLM pattern-matched to "95th percentile= -1.645" from normal tables without understanding domain semantics\citep{kambhampati2024llms}.

\subsubsection{Case Study 3: Library API Misuse (Expected Shortfall)}

Expected Shortfall (also called Conditional VaR) provides a more comprehensive risk measure than VaR \citep{embrechts1997modelling,rockafellar2000optimization}.

\paragraph{Ground Truth:}
\begin{equation}
\text{ES}_{\alpha} = -\mu + \sigma \cdot \frac{\phi(\Phi^{-1}(\alpha))}{\alpha}
\end{equation}

\paragraph{LLM Output:}
\begin{verbatim}
from scipy.stats import norm

def formula(portfolio_value, daily_volatility, alpha):
    z_alpha = norm.ppf(alpha)
    phi_z = norm.pdf(z_alpha)  # Returns frozen distribution!
    return portfolio_value * daily_volatility * phi_z / alpha
    # ERROR: Can't multiply float by rv_continuous_frozen
\end{verbatim}

\paragraph{Analysis:}
\begin{itemize}
\item Incorrect scipy API usage (\texttt{norm.pdf(x)} returns frozen object, not scalar)
\item Runtime error makes failure detectable (unlike Case Study 1)
\item LLMs struggle with specialized library APIs having non-obvious calling conventions\citep{kambhampati2024llms}
\item Correct usage requires \texttt{norm.pdf(z\_alpha, 0, 1)} or frozen distribution methods
\end{itemize}

\subsection{Success Pattern Analysis}

Table~\ref{tab:llm_success_patterns} shows LLM performance by formula complexity\citep{kambhampati2024llms}.

\begin{table}[htbp]
\centering
\caption{LLM Success Rate by Formula Complexity}
\label{tab:llm_success_patterns}
\begin{tabular}{@{}lccp{5cm}@{}}
\toprule
\textbf{Complexity} & \textbf{Success} & \textbf{Avg $R^2$} & \textbf{Examples} \\
\midrule
Linear (2 vars) & 100\% & 1.000 & $\sigma = E\varepsilon$, $V = IR$ \\
Polynomial (2-3 vars) & 100\% & 1.000 & $E = \frac{1}{2}mv^2$, $F = ma$ \\
Rational (3-4 vars) & 95\% & 0.995 & $Y = k/X$, Re $= vL/\nu$ \\
Exponential & 100\% & 1.000 & $k = A \exp(-E_a/RT)$ \\
Square root & 60\% & 0.600 & $\text{IL} = 2\sqrt{r}/(1+r) - 1$ \\
Statistical (scipy) & 20\% & $-2.1$ & VaR, ES with \texttt{norm.ppf/pdf} \\
Multi-step & 0\% & --- & Liquidation, portfolio optimization \\
\bottomrule
\end{tabular}
\end{table}

\paragraph{Training Data Hypothesis:}

\begin{table}[htbp]
\centering
\caption{LLM Success vs. Estimated Training Exposure}
\label{tab:training_exposure}
\begin{tabular}{@{}lcc@{}}
\toprule
\textbf{Domain} & \textbf{Training Likelihood} & \textbf{Success Rate} \\
\midrule
Classical Physics & Extensive (textbooks, Wikipedia) & 100\% \\
Engineering Mechanics & Extensive (handbooks, standards) & 100\% \\
Basic Chemistry & Extensive (textbooks, databases) & 75\% \\
DeFi/AMM & Limited (recent, niche, evolving) & 65\% \\
Risk Analytics & Limited (specialized, proprietary) & 40\% \\
\bottomrule
\end{tabular}
\end{table}

\subsection{Extrapolation Testing}

We evaluated formulas on data 2$\times$ outside the training distribution (e.g., if trained on $x \in [1, 10]$, tested on $x \in [20, 40]$).

\begin{table}[htbp]
\centering
\caption{Extrapolation Performance}
\label{tab:extrapolation}
\begin{tabular}{@{}lcc@{}}
\toprule
\textbf{Method} & \textbf{In-Distribution $R^2$} & \textbf{Extrapolation Error} \\
\midrule
Pure LLM (GPT-4) & 0\citep{bubeck2023sparks}\citep{romera2023mathematical}\citep{kambhampati2024llms}.89 & 847\% \\
Pure LLM (Claude Opus 4) & 0\citep{kambhampati2024llms}.93 & 412\% \\
\midrule
Symbolic (PySR) & 0.94 & 23\% \\
HypatiaX (DeFi) & 0.999 & 8\% \\
\bottomrule
\end{tabular}
\end{table}

\paragraph{Analysis:} LLM-generated formulas often overfit to training data patterns, producing nonsensical predictions outside the training regime (e.g., negative probabilities, domain violations). This extrapolation failure is well-documented in neural network literature \citep{zhang2024extrapolation}, where models exhibit spectral bias toward low-frequency functions. Symbolic methods enforce physical constraints that ensure extrapolation validity.

\subsection{Implications}

These empirical results demonstrate:

\begin{enumerate}
\item \textbf{LLMs are reproductive, not generative}: 100\% success on classical formulas versus 45\% on novel domains confirms reproduction of training patterns\citep{kambhampati2024llms}.

\item \textbf{Silent failures are dangerous}: High $R^2$ in training range doesn't imply correctness. Unit errors and sign flips produce plausible but catastrophically wrong results.

\item \textbf{Prompt engineering cannot overcome structural limitations}: We tested multiple prompt strategies with no improvement on novel domains.

\item \textbf{Multi-layer validation is essential}: Dimensional analysis, domain constraints, and extrapolation testing detect failures missed by $R^2$ alone.
\end{enumerate}

These findings motivate our hybrid architecture, where symbolic methods perform discovery under physics constraints while LLMs provide initialization and interpretation\citep{koza1994genetic}\citep{kambhampati2024llms}.

\section{Theoretical Framework}
\label{sec:theory}

\subsection{Formal Definitions}

\begin{definition}[Analytical Expression Discovery]
\label{def:analytical_discovery}
Let $\mathcal{D} = \{(x_i, y_i)\}_{i=1}^n$ where $x_i \in \R^d$, $y_i \in \R$\citep{koza1994genetic}. An analytical expression discovery system produces $E: \R^d \to \R$ satisfying:
\begin{enumerate}
\item \textbf{Symbolic form}: $E$ is expressible using operators $\mathcal{O}$ in closed form
\item \textbf{Physical consistency}: $E$ satisfies domain constraints $\mathcal{C}$
\item \textbf{Empirical accuracy}: Loss $\mathcal{L}(E, \mathcal{D}) < \epsilon$
\item \textbf{Extrapolation}: $E$ generalizes beyond training regime
\item \textbf{Parsimony}: $E$ has minimal complexity given above constraints
\end{enumerate}
\end{definition}

This definition draws on established principles of symbolic computation \citep{meurer2017sympy} and dimensional analysis \citep{buckingham1914physically}.

\begin{definition}[Reproductive System]
\label{def:reproductive}
System $S$ is \emph{reproductive} if generation probability for expressions in training data vastly exceeds that for novel expressions:
\[
P_S(E_{\text{out}} = E \mid E \in \mathcal{C}_{\text{train}}) \gg P_S(E_{\text{out}} = E \mid E \notin \mathcal{C}_{\text{train}})
\]
where $\mathcal{C}_{\text{train}}$ is the training corpus.
\end{definition}

\subsection{Main Theoretical Result}

\begin{theorem}[LLMs are Reproductive Systems]
\label{thm:llm_reproductive}
Autoregressive language models are reproductive systems in the sense of Definition~\ref{def:reproductive}\citep{lewkowycz2022minerva}.
\end{theorem}

\begin{proof}[Proof Sketch]
Expression $E = t_1 \cdots t_n$ is generated token-by-token via:
\[
P(E \mid D) = \prod_{i=1}^n P(t_i \mid t_1, \ldots, t_{i-1}, D; \theta)
\]

Training optimizes $\theta$ to maximize log-likelihood on corpus $\mathcal{C}$:
\[
\theta^* = \arg\max_\theta \sum_{E \in \mathcal{C}} \log P(E; \theta)
\]

This concentrates probability mass on frequently observed token sequences. For expression $E \notin \mathcal{C}$, let $d(E, \mathcal{C})$ denote minimum token-level edit distance to nearest training expression. Each unusual token incurs probability penalty. Approximating local independence of penalties:
\[
P(E \mid D) \lesssim \prod_{i \in \text{unusual}(E)} p_{\min} \leq p_{\min}^{d(E, \mathcal{C})}
\]
where $p_{\min}$ is baseline token probability. Since $p_{\min} < 1$, this decays exponentially:
\[
P(E \mid D) = O(\exp(-\beta \cdot d(E, \mathcal{C})))
\]
for some $\beta > 0$ depending on model temperature and vocabulary size.

Thus, for expressions requiring $d > k$ novel tokens, $P(E \mid D) \to 0$ exponentially fast, establishing reproductive behavior. \qed
\end{proof}

This reproductive behavior is consistent with observations about large language model training \citep{hoffmann2022training} and theoretical analysis of fundamental scaling limits \citep{xu2025fundamental}.

\paragraph{Remark:} This proof assumes local independence of token probabilities, which is approximate. A rigorous proof would require analyzing the full joint distribution, but empirical evidence (Section~\ref{sec:llm_limitations}) strongly supports the conclusion\citep{kambhampati2024llms}.

\subsection{Implications for Discovery}\citep{koza1994genetic}

Theorem~\ref{thm:llm_reproductive} has several implications:\citep{kambhampati2024llms}

\begin{enumerate}
\item \textbf{Training coverage determines capability}: LLMs can only reliably generate expressions within small edit distance of training data\citep{kambhampati2024llms}.

\item \textbf{Novel domains require exponentially unlikely generation}: Specialized formulas (DeFi, advanced statistics) requiring novel token combinations have vanishingly small generation probability.

\item \textbf{Scale alone cannot overcome this limitation}: Increasing model size or training data improves coverage but doesn't enable genuine extrapolation beyond the training distribution.

\item \textbf{Hybrid approaches are necessary}: Combining LLM pattern recognition with symbolic search enables both efficiency (warm-starting) and genuine discovery\citep{koza1994genetic}\citep{kambhampati2024llms}.
\end{enumerate}

\section{HypatiaX Architecture}
\label{sec:architecture}

Motivated by empirical limitations (Section~\ref{sec:llm_limitations}) and theoretical analysis (Section~\ref{sec:theory}), we introduce \textsc{HypatiaX} (Hybrid Physics-Aware Analytical eXpression discovery), a system where symbolic methods perform mathematical reasoning while LLMs provide natural language interpretation\citep{koza1994genetic}\citep{lewkowycz2022minerva}\citep{kambhampati2024llms}.

\subsection{Design Principles}

\begin{enumerate}
\item \textbf{Symbolic methods for discovery}: Evolutionary search with physics constraints finds mathematically valid expressions\citep{koza1994genetic}
\item \textbf{Multi-layer validation}: Dimensional analysis, domain constraints, execution verification, and extrapolation testing
\item \textbf{LLMs as interfaces}: Natural language interpretation of results and optional warm-starting of search\citep{kambhampati2024llms}
\item \textbf{Failure transparency}: Invalid expressions explicitly rejected, not silently accepted
\end{enumerate}

\subsection{System Architecture}

\textsc{HypatiaX} integrates five layers:

\begin{figure}[htbp]
\centering
\begin{small}
\begin{verbatim}
+-----------------------------------------------+
|      MULTIMODAL DATA SOURCES                  |
|  [Images]  [Time-Series]  [Text]  [Tables]    |
+-----------------------------------------------+
                 |
+-----------------------------------------------+
|      PREPROCESSING LAYER                      |
|  * Feature Extraction                         |
|  * Data Normalization                         |
|  * Variable Identification                    |
+-----------------------------------------------+
                 |
+-----------------------------------------------+
|   SYMBOLIC DISCOVERY CORE (PySR)              |\citep{koza1994genetic}
|   * Multi-Objective Evolutionary Search       |
|   * Physics-Informed Operators                |
|   * Constraint Satisfaction                   |
|   * Pareto Front Optimization                 |
+-----------------------------------------------+
                 |
+-----------------------------------------------+
|      VALIDATION FRAMEWORK                     |
|  * Dimensional Analysis                       |
|  * Domain Constraint Checking                 |
|  * Numerical Stability Testing                |
|  * Extrapolation Verification                 |
+-----------------------------------------------+
                 |
+-----------------------------------------------+
|      LLM INTERPRETATION LAYER (Optional)      |\citep{kambhampati2024llms}
|  * Physical Explanation Generation            |
|  * Literature Contextualization               |
|  * Natural Language Reporting                 |
+-----------------------------------------------+
                 |
+-----------------------------------------------+
|      OUTPUT: Validated Analytical Report      |
+-----------------------------------------------+
\end{verbatim}

\end{small}
\caption{HypatiaX hybrid symbolic-neural architecture. Symbolic methods perform discovery; LLMs interpret results.}
\label{fig:architecture}
\end{figure}

\subsection{Symbolic Discovery Core}\citep{koza1994genetic}

We employ PySR \citep{cranmer2023interpretable}, building on genetic programming principles \citep{koza1992genetic,fortin2012deap}, for multi-objective genetic programming:

\begin{align}
\text{Minimize:} \quad & \mathcal{L}(E, \mathcal{D}) = \frac{1}{n}\sum_{i=1}^n (E(x_i) - y_i)^2 \label{eq:loss}\\
\text{Minimize:} \quad & \text{Complexity}(E) = \sum_{t \in E} w(t) \label{eq:complexity}\\
\text{Subject to:} \quad & E \in \mathcal{C} \label{eq:constraints}
\end{align}

where $w(t)$ assigns complexity weights to operators and $\mathcal{C}$ encodes physical constraints.

\paragraph{Operator Set:} $\mathcal{O} = \{+, -, \times, \div, \exp, \log, \sin, \cos, x^n, \sqrt{x}\}$

\paragraph{Evolutionary Operators:}
\begin{itemize}
\item \textbf{Crossover}: Subtree exchange between expressions
\item \textbf{Mutation}: Random operator/constant modification
\item \textbf{Simplification}: Symbolic reduction using SymPy
\item \textbf{Selection}: Tournament selection on Pareto front
\end{itemize}

\subsection{Multi-Layer Validation Framework}

\subsubsection{Layer 1: Dimensional Analysis}

Every physical quantity has dimensions $[M^a L^b T^c \cdots]$. Valid expressions satisfy:

\begin{proposition}[Dimensional Homogeneity]
\label{prop:dimensional}
For expression $E$ combining quantities $q_1, \ldots, q_k$ with dimensions $[d_1], \ldots, [d_k]$, if $E$ is physically valid, then all terms added/subtracted must have identical dimensions.
\end{proposition}

\textbf{Implementation:} We track dimensions symbolically and reject expressions violating homogeneity.

\subsubsection{Layer 2: Domain Constraint Checking}

Domain-specific rules encode physical laws:

\begin{itemize}
\item \textbf{Conservation laws}: Energy, momentum, mass conservation
\item \textbf{Thermodynamic constraints}: Entropy non-decrease, positive absolute temperature
\item \textbf{Financial no-arbitrage}: Portfolio values, probability constraints
\item \textbf{Statistical validity}: Probabilities in $[0,1]$, valid distribution parameters
\end{itemize}

\subsubsection{Layer 3: Numerical Stability Testing}

Expressions must produce finite values over admissible domains. We evaluate on dense grids and reject if:

\begin{equation}
\frac{|\{i : |E(x_i)| < \infty\}|}{N} < \alpha \quad (\text{default } \alpha = 0.95)
\end{equation}

\subsubsection{Layer 4: Extrapolation Verification}

We partition data into training $\mathcal{D}_{\text{train}}$ and extrapolation $\mathcal{D}_{\text{extrap}}$ regions separated by 2$\times$ in feature space. Valid expressions must satisfy:

\begin{equation}
R^2(\mathcal{D}_{\text{extrap}}) > \tau_{\text{extrap}} \quad (\text{default } \tau_{\text{extrap}} = 0.80)
\end{equation}

\subsection{Domain-Aware Singularity Reconciliation}

Symbolic analysis may flag potential singularities (e.g., division by zero) that are unreachable under domain constraints.

\begin{proposition}[Safe Singularities]
\label{prop:safe_singularity}
Let $E(x) = g(x)/h(x)$. If symbolic analysis detects $\exists x_0: h(x_0) = 0$ but domain constraints $\mathcal{C}$ ensure $h(x) \neq 0$ for all admissible $x \in \mathcal{C}$, then the singularity is safe and $E$ is accepted.
\end{proposition}

\textbf{Example (Michaelis-Menten):} The Michaelis-Menten equation \citep{michaelis1913kinetics}, $v(S) = \frac{V_{\max}S}{K_m + S}$, has potential singularity at $K_m + S = 0$. Domain constraints $K_m > 0, S > 0$ ensure $K_m + S > 0$ always, so the expression is accepted.

\subsection{Acceptance Criteria}

Expressions are accepted if they satisfy:

\begin{equation}
\text{Accept} \iff \begin{cases}
R^2 > 0.99 \land V > 30 & \text{(High accuracy)} \\
R^2 > 0.95 \land V > 80 & \text{(Strong validation)}
\end{cases}
\end{equation}

where $V \in [0, 100]$ is aggregate validation score across all layers.

\section{LLM-Enhanced Variant}\citep{kambhampati2024llms}
\label{sec:llm_enhancement}

While Section~\ref{sec:llm_limitations} established that pure LLMs cannot reliably perform discovery, they excel at pattern recognition for \emph{known} formulas. We leverage this via optional warm-starting\citep{koza1994genetic}\citep{kambhampati2024llms}.

\subsection{Hybrid Discovery Mode}\citep{koza1994genetic}

\begin{algorithm}[htbp]
\caption{LLM-Guided Symbolic Discovery}
\label{alg:hybrid}
\begin{algorithmic}[1]
\STATE \textbf{Input:} Data $\mathcal{D}$, domain context, variables
\STATE \textbf{Output:} Validated expression $E^*$
\STATE
\STATE $E_{\text{LLM}} \gets$ LLM\citep{kambhampati2024llms}.generate($\mathcal{D}$, context)
\STATE $R^2_{\text{LLM}} \gets$ evaluate($E_{\text{LLM}}$, $\mathcal{D}$)\citep{kambhampati2024llms}
\IF{$R^2_{\text{LLM}} \geq 0\citep{kambhampati2024llms}.95$ \AND validate($E_{\text{LLM}}$) = PASS}
    \STATE \textbf{return} $E_{\text{LLM}}$ \COMMENT{Accept LLM solution}\citep{kambhampati2024llms}
\ELSE
    \STATE $P_0 \gets$ initialize\_population($E_{\text{LLM}}$)\citep{kambhampati2024llms}
    \STATE $E^* \gets$ PySR.evolve($P_0$, $\mathcal{D}$, $\mathcal{C}$)
    \STATE \textbf{return} $E^*$ \COMMENT{Refined via symbolic search}
\ENDIF
\end{algorithmic}
\end{algorithm}

\subsection{LLM Prompting Strategy}\citep{kambhampati2024llms}

We employ structured prompts containing:

\begin{small}
\begin{verbatim}
You are a mathematical physicist discovering equations from data.

Domain: {domain_name}
Variables:
  - {var1}: {description}, units: {units}, range: [{min}, {max}]
  - {var2}: ...

Data characteristics:
  - Correlations: {correlation_matrix}
  - Detected patterns: {linear/exponential/...}
  - Sample size: {n}

Generate:
1. Top 3 candidate expressions (LaTeX format)
2. Confidence score (0-1) for each
3. Physical interpretation
4. Expected asymptotic behavior

Constraints:
- Dimensionally consistent
- Use only: +, -, *, /, ^, exp, log, sin, cos
- No undefined functions
- Parsable by SymPy
\end{verbatim}
\end{small}

\paragraph{Critical Design Choices:}

\begin{itemize}
\item \textbf{No code generation}: LLM produces mathematical expressions only, avoiding syntax errors\citep{kambhampati2024llms}
\item \textbf{Domain context}: Semantic variable descriptions enable domain knowledge application
\item \textbf{Statistical hints}: Correlation patterns guide toward likely functional forms
\item \textbf{Multiple candidates}: Requesting 3 alternatives allows fallback if top fails
\item \textbf{Strict constraints}: SymPy-compatible operators ensure evaluability
\end{itemize}

\subsection{Performance Analysis}

Table~\ref{tab:llm_hybrid_performance} compares pure symbolic versus LLM-guided discovery\citep{koza1994genetic}\citep{kambhampati2024llms}.

\begin{table}[htbp]
\centering
\caption{LLM-Guided Discovery Performance (30 tests)}
\label{tab:llm_hybrid_performance}
\begin{tabular}{@{}lccccc@{}}
\toprule
\textbf{Method} & \textbf{Success} & \textbf{Avg Time} & \textbf{Avg $R^2$} & \textbf{Speedup} & \textbf{Cost} \\
\midrule
Pure PySR & 88.9\% & 390s & 0.9987 & 1.0$\times$ & \$0 \\
LLM-only & 73.0\% & 13.5s & 1.0000 & 28.9$\times$ & \$0\citep{kambhampati2024llms}.02 \\
\textbf{LLM + PySR (hybrid)} & \textbf{96.7\%} & \textbf{70.0s} & \textbf{1.0000} & \textbf{5.6}$\times$ & \$0\citep{kambhampati2024llms}.02 \\
\bottomrule
\end{tabular}
\end{table}

\paragraph{Discovery Mode Breakdown:}\citep{koza1994genetic}

\begin{table}[htbp]
\centering
\caption{Discovery Modes in LLM-Guided System (30 tests)}
\label{tab:discovery_modes}
\begin{tabular}{@{}lcccc@{}}
\toprule
\textbf{Mode} & \textbf{Tests} & \textbf{Success} & \textbf{Avg Time} & \textbf{Avg $R^2$} \\
\midrule
LLM-only accepted & 22 & 100\% & 13.5s & 1\citep{kambhampati2024llms}.0000 \\
LLM + PySR refinement & 7 & 85.7\% & 226.8s & 0\citep{kambhampati2024llms}.9998 \\
PySR fallback & 1 & 0\% & 360.3s & --- \\
\midrule
\textbf{Combined} & \textbf{30} & \textbf{96.7\%} & \textbf{70.0s} & \textbf{1.0000} \\
\bottomrule
\end{tabular}
\end{table}

\paragraph{Key Findings:}

\begin{enumerate}
\item \textbf{82\% time reduction}: Average 70s versus 390s for pure PySR
\item \textbf{73\% LLM-only success}: For simple well-known relationships, LLM initialization suffices\citep{kambhampati2024llms}
\item \textbf{Success rate improvement}: 96.7\% versus 88.9\% for pure PySR (though not statistically significant with $p = 0.18$, $n = 30$)
\item \textbf{Preserved guarantees}: All LLM outputs undergo identical validation as PySR expressions\citep{kambhampati2024llms}
\item \textbf{Minimal cost}: \$0.02 per test, negligible compared to computational savings
\end{enumerate}

\subsection{When LLM Guidance Succeeds}\citep{kambhampati2024llms}

\paragraph{LLM-only acceptance (73\% of tests):}\citep{kambhampati2024llms}
\begin{itemize}
\item Linear relationships: Ohm's law, Hooke's law, ideal gas law
\item Power laws: Kinetic energy, gravitational force
\item Exponential decay: Cooling, radioactive decay, Arrhenius
\item Common DeFi: Constant product AMM, simple fee calculations
\end{itemize}

\paragraph{LLM + PySR refinement (23\% of tests):}\citep{kambhampati2024llms}
\begin{itemize}
\item Multi-term expressions: Stefan-Boltzmann, Bernoulli equation
\item Coupled interactions: Capital efficiency, portfolio risk with correlations
\item Non-standard forms: Kelly criterion \citep{kelly1956criterion}, impermanent loss
\end{itemize}

\paragraph{Complete failure (4\% of tests):}
\begin{itemize}
\item Hooke's Law confusion: LLM generated elastic potential energy $U = \frac{1}{2}kx^2$ instead of force $F = kx$\citep{kambhampati2024llms}
\item PySR overfit to noise, validation correctly rejected (score 12.1/100)
\end{itemize}

\subsection{Design Rationale}

The LLM-guided architecture succeeds by exploiting complementary strengths:\citep{kambhampati2024llms}

\begin{itemize}
\item \textbf{LLM strengths}: Pattern recognition for known formulas, domain knowledge application, human-readable output, speed (10-15s)\citep{kambhampati2024llms}
\item \textbf{Symbolic strengths}: Exhaustive search, novel discovery, dimensional consistency, reliable extrapolation, formal guarantees\citep{koza1994genetic}
\end{itemize}

Using LLM initialization as warm-start rather than replacement achieves both efficiency (82\% time reduction) and reliability (all validation layers apply)\citep{kambhampati2024llms}.

%Phase 2: Insert Experimental Validation (File 4 → Section 7)
%Part 1: Section Header and Problem Formulation

\section{Experimental Validation}
\label{sec:experiments}

We present comprehensive experimental results across 15 ground truth equations spanning five scientific and financial domains, plus extended validation across 131 total tests. Our evaluation compares pure LLM formula generation, neural network regression, pure symbolic regression (PySR), and our hybrid HypatiaX system\citep{la2021contemporary}\citep{virtanen2020scipy}\citep{kambhampati2024llms}.

\subsection{Problem Formulation}

\begin{definition}[Symbolic Discovery Problem]
Given training data $\mathcal{D}_{\text{train}} = \{(\mathbf{x}_i, y_i)\}_{i=1}^N$ where $\mathbf{x}_i \in \mathbb{R}^d$ and $y_i \in \mathbb{R}$, discover a mathematical expression $f: \mathbb{R}^d \to \mathbb{R}$ such that:
\begin{enumerate}
    \item \textbf{Accuracy}: $\text{RMSE}_{\text{train}}(f) = \sqrt{\frac{1}{N}\sum_{i=1}^N (f(\mathbf{x}_i) - y_i)^2} < \epsilon$
    \item \textbf{Parsimony}: Complexity $C(f)$ is minimal among equivalent expressions
    \item \textbf{Extrapolation}: $\text{RMSE}_{\text{extrap}}(f) \approx \text{RMSE}_{\text{train}}(f)$ for $\mathbf{x} \notin \text{Conv}(\mathcal{D}_{\text{train}})$
\end{enumerate}
\end{definition}

Traditional machine learning optimizes only criterion (1). We argue that criterion (3)---extrapolation---is the litmus test for discovering true functional relationships versus memorizing patterns\citep{molnar2020interpretable}\citep{rudin2019stop}.

%Phase 2: Part 2 - Experimental Design
\subsection{Experimental Design}

\subsubsection{Ground Truth Test Suite}

We constructed a core test suite of 15 equations with known analytical forms (Table~\ref{tab:ground_truth_suite}), selected to satisfy:

\begin{enumerate}
    \item \textbf{Analytical form}: Closed-form mathematical expression exists
    \item \textbf{Physical grounding}: Derived from first principles or established empirical laws
    \item \textbf{Practical relevance}: Used in real-world scientific or financial applications\citep{virtanen2020scipy}
    \item \textbf{Variable diversity}: Tests handling different dimensionalities (1-4 variables)
\end{enumerate}

\begin{table}[htbp]
\centering
\caption{Ground Truth Test Suite: 15 Core Equations Across 5 Domains}
\label{tab:ground_truth_suite}
\small
\begin{tabular}{llcc}
\toprule
\textbf{Domain} & \textbf{Equation} & \textbf{Variables} & \textbf{Form} \\
\midrule
\multirow{3}{*}{Chemistry}
& Arrhenius & 1 & $k = A\exp(-E_a/RT)$ \\
& Henderson-Hasselbalch & 2 & $\text{pH} = \text{p}K_a + \log_{10}([A^-]/[\text{HA}])$ \\
& Rate Law & 2 & $r = k[A]^2[B]$ \\
\midrule
\multirow{3}{*}{Biology}
& Allometric Scaling & 1 & $Y = aM^b$ \\
& Michaelis-Menten & 1 & $v = V_{\max}[S]/(K_m + [S])$ \\
& Logistic Growth & 1 & $dN/dt = rN(1-N/K)$ \\
\midrule
\multirow{3}{*}{Physics}
& Kinetic Energy & 2 & $E = \frac{1}{2}mv^2$ \\
& Gravitational Force & 3 & $F = Gm_1m_2/r^2$ \\
& Ideal Gas Law & 3 & $P = nRT/V$ \\
\midrule
\multirow{3}{*}{DeFi AMM}
& Impermanent Loss & 1 & $\text{IL} = 2\sqrt{r}/(1+r) - 1$ \\
& Price Impact & 2 & $\Delta p = \Delta x/(x + \Delta x)$ \\
& Constant Product & 1 & $y = k/x$ \\
\midrule
\multirow{3}{*}{DeFi Risk}
& Value-at-Risk 95\% & 2 & $\text{VaR} = P\sigma \cdot 1.645$ \\
& Liquidation Price & 2 & $p_{\text{liq}} = p_0(1 - 1/(L \cdot m))$ \\
& Portfolio Variance & 3 & $\sigma_p = \sqrt{\sigma_1^2 + \sigma_2^2 + 2\rho\sigma_1\sigma_2}$ \\
\bottomrule
\end{tabular}
\end{table}

\paragraph{Data Generation Protocol}

For each equation:
\begin{itemize}
    \item \textbf{Sample size}: $N = 200$ observations (80\% train, 20\% test)
    \item \textbf{Noise model}: $y_i = f(\mathbf{x}_i) + \epsilon_i$ where $\epsilon_i \sim \mathcal{N}(0, \sigma^2)$, $\sigma = 0.05 \cdot \text{std}(f(\mathbf{x}))$
    \item \textbf{Variable ranges}: Physically/financially meaningful domains
    \item \textbf{Random seed}: Fixed at 42 for reproducibility
\end{itemize}

\subsubsection{Extended Test Suite}

Beyond the 15 core equations, we conducted extended validation across:
\begin{itemize}
\item \textbf{Pure LLM Baseline}: 40 tests (20 classical physics, 20 DeFi)\citep{kambhampati2024llms}
\item \textbf{Pure Symbolic (PySR)}: 18 tests across multiple domains
\item \textbf{LLM-Guided Hybrid}: 30 tests evaluating warm-starting benefits\citep{kambhampati2024llms}
\item \textbf{DeFi-Optimized}: 23 tests with domain-specific tuning
\item \textbf{Total}: 131 distinct test cases
\end{itemize}

%Phase 2: Part 3 - Extrapolation Test Protocol

\subsection{Extrapolation Test Protocol}
\label{sec:extrapolation_metric}

\subsubsection{Test Set Generation}

We define extrapolation capability as maintaining accuracy beyond the training domain. Given training data $\mathcal{D}_{\text{train}}$ from range $[\mathbf{x}_{\min}, \mathbf{x}_{\max}]$, we generate test data $\mathcal{D}_{\text{extrap}}$ at three distances:

\begin{definition}[Extrapolation Regimes]
Let $\alpha$ be the extrapolation factor:
\begin{align}
\text{Near (1.2\times):} \quad &\mathbf{x} \in [1.2\mathbf{x}_{\max}, 1.44\mathbf{x}_{\max}] \\
\text{Medium (2\times):} \quad &\mathbf{x} \in [2.0\mathbf{x}_{\max}, 3.0\mathbf{x}_{\max}] \\
\text{Far (5\times):} \quad &\mathbf{x} \in [5.0\mathbf{x}_{\max}, 7.5\mathbf{x}_{\max}]
\end{align}
\end{definition}

For multivariate functions, all variables are extrapolated simultaneously by the same factor $\alpha$ to test behavior in truly novel regions of the input space.

\subsubsection{Extrapolation Error Metric}

The extrapolation error metric quantifies degradation:

\begin{definition}[Extrapolation Error]
\label{def:extrap_error}
Given a learned model $\hat{f}$ with training RMSE $\text{RMSE}_{\text{train}}$ and extrapolation test set $\mathcal{D}_{\text{extrap}}$:
\begin{equation}
E_{\text{extrap}} = \frac{\text{RMSE}(\hat{f}, \mathcal{D}_{\text{extrap}})}{\text{RMSE}(\hat{f}, \mathcal{D}_{\text{train}})} \times 100\%
\label{eq:extrap_error}
\end{equation}
\end{definition}

\textbf{Interpretation}:
\begin{itemize}
    \item $E_{\text{extrap}} = 100\%$: Model maintains training accuracy
    \item $E_{\text{extrap}} = 0\%$: Perfect extrapolation (exact functional form)
    \item $E_{\text{extrap}} > 1000\%$: Catastrophic failure
\end{itemize}

\begin{example}[Arrhenius Equation Extrapolation]
For the Arrhenius equation $k = A \exp(-E_a/RT)$ with temperature $T$:
\begin{itemize}
    \item Training: $T \in [273, 373]$K
    \item Medium extrapolation: $T \in [746, 1119]$K ($2\times$ max)
    \item Neural Network: $\text{RMSE}_{\text{train}} = 0.05$, $\text{RMSE}_{\text{extrap}} = 167.4$
    \item $E_{\text{extrap}} = (167.4 / 0.05) \times 100\% = 3348\%$
\end{itemize}
This indicates the neural network's predictions are $33\times$ worse than training error.
\end{example}

\subsubsection{Statistical Analysis}

We assessed significance using:

\begin{itemize}
    \item \textbf{Mann-Whitney U test}: Non-parametric comparison of Hybrid v40 vs Neural Network extrapolation errors. Null hypothesis: $H_0: E_{\text{hybrid}} \geq E_{\text{NN}}$. One-tailed test with $\alpha = 0.05$.

    \item \textbf{Effect size (Cohen's d)}: Quantifies practical significance:
    \begin{equation}
    d = \frac{\mu_{\text{NN}} - \mu_{\text{hybrid}}}{\sqrt{(\sigma_{\text{NN}}^2 + \sigma_{\text{hybrid}}^2)/2}}
    \end{equation}
    Interpretation: $d > 0.8$ (large effect), $d > 2.0$ (huge effect).

    \item \textbf{95\% confidence intervals}: Calculated using Welch's t-distribution for unequal variances.
\end{itemize}

%Phase 2: Part 4 - Methods Under Evaluation (Method 1: Pure LLM)

\subsection{Methods Under Evaluation}

\subsubsection{Method 1: Pure LLM Baseline}\citep{kambhampati2024llms}

We implemented a baseline using Claude Sonnet 4 with specialized prompt engineering:

\begin{lstlisting}[language=Python, caption=Pure LLM Formula Generation]
class PureLLMBaseline:
    def generate_formula(self, description, domain,
                        variable_names, metadata):
        prompt = f"""You are a scientific equation discovery system\citep{koza1994genetic}\citep{virtanen2020scipy}.

Domain: {domain}
Task: {description}
Variables: {variable_names}
Units: {metadata['units']}

Generate the mathematical relationship as:
y = f(x1, x2, ..., xn)

Provide ONLY the formula, no explanation."""

        response = anthropic.messages.create(
            model="claude-sonnet-4-20250514",
            max_tokens=1024,
            messages=[{"role": "user", "content": prompt}]
        )

        return self.parse_and_compile(response.content[0].text)

    def parse_and_compile(self, formula_str):
        """Convert LaTeX/Unicode to executable Python function"""
        import sympy as sp
        expr = sp.sympify(formula_str)
        func = sp.lambdify(symbols, expr, modules=['numpy'])
        return {"formula": formula_str, "function": func}
\end{lstlisting}

\textbf{Strengths}: Fast execution ($\sim$7s per test), generates human-readable formulas, perfect interpolation ($R^2 = 1.0$).

\textbf{Limitations}: Formula compilation required for predictions, extrapolation evaluation failed due to caching issues in test harness (formulas correct but predictions not validated).

%Phase 2: Part 5 - Method 2: Neural Network Baseline
\subsubsection{Method 2: Neural Network Baseline}

Standard multilayer perceptron (MLP) implemented in PyTorch:

\begin{lstlisting}[language=Python, caption=Neural Network Architecture]
class SimpleNN(nn.Module):
    def __init__(self, input_dim):
        super().__init__()
        self.network = nn.Sequential(
            nn.Linear(input_dim, 64),
            nn.ReLU(),
            nn.Linear(64, 32),
            nn.ReLU(),
            nn.Linear(32, 1)
        )

    def forward(self, x):
        return self.network(x)

# Training procedure
def train_and_evaluate(X, y, epochs=200):
    scaler_X = StandardScaler()
    scaler_y = StandardScaler()

    X_scaled = scaler_X.fit_transform(X)
    y_scaled = scaler_y.fit_transform(y.reshape(-1, 1)).flatten()

    model = SimpleNN(X.shape[1])
    optimizer = torch.optim.Adam(model.parameters(), lr=0.01)
    criterion = nn.MSELoss()

    X_train, X_test, y_train, y_test = train_test_split(
        X_scaled, y_scaled, test_size=0.2, random_state=42
    )

    # Training loop (200 epochs)...

    return {
        "model": model,
        "scaler_X": scaler_X,
        "scaler_y": scaler_y,
        "r2": r2_score,
        "rmse": rmse
    }
\end{lstlisting}

Critical for extrapolation: we store both the trained model and normalization scalers to enable proper prediction on out-of-distribution data:

\begin{lstlisting}[language=Python, caption=Neural Network Prediction with Proper Scaling]
def nn_predict(model, scaler_X, scaler_y, X_new):
    """Make predictions on new data with proper scaling"""
    model.eval()

    # Scale input using training statistics
    X_scaled = scaler_X.transform(X_new)
    X_tensor = torch.FloatTensor(X_scaled)

    # Predict in scaled space
    with torch.no_grad():
        y_scaled = model(X_tensor).numpy().flatten()

    # Transform back to original scale
    y_pred = scaler_y.inverse_transform(
        y_scaled.reshape(-1, 1)
    ).flatten()

    return y_pred
\end{lstlisting}

\textbf{Strengths}: Fastest method ($\sim$2s per test), reliable interpolation.

\textbf{Limitations}: Black-box predictions, no interpretable formula, catastrophic extrapolation failure (quantified in Section~\ref{sec:results}).

%Phase 2: Part 6 - Method 3: Hybrid System v40 (Stage 1-2)

\subsubsection{Method 3: Hybrid System v40 (Our Method)}

Our hybrid approach integrates LLM guidance with symbolic regression through four stages:\citep{la2021contemporary}\citep{kambhampati2024llms}

\paragraph{Stage 1: LLM-Guided Hypothesis Generation}\citep{kambhampati2024llms}

We use Claude Sonnet 4 to generate candidate symbolic expressions incorporating domain knowledge:

\begin{lstlisting}[language=Python, caption=LLM Hypothesis Generation]
llm_config = LLMConfig(
    enabled=True,
    api_key=ANTHROPIC_API_KEY,
    n_candidates=5,
    model="claude-sonnet-4-20250514"
)

# Generate candidate expressions
candidates = llm\citep{kambhampati2024llms}.generate_symbolic_hypotheses(
    description=description,
    domain=domain,
    variable_names=variable_names,
    variable_units=units,
    n_candidates=5
)
# Example output for Kinetic Energy:
# ["0.5 * m * v^2", "m * v^2", "sqrt(m) * v", ...]
\end{lstlisting}

\paragraph{Stage 2: Variable Name Sanitization}

Julia/PySR reserves single-letter variable names (S, N, C, D, E, I, O) for internal use. We implement automatic sanitization:

\begin{lstlisting}[language=Python, caption=Variable Name Sanitization]
class VariableNameSanitizer:
    RESERVED_NAMES = {'S', 'N', 'C', 'D', 'E', 'I', 'O'}

    def sanitize(self, variable_names):
        sanitized = []
        mapping = {}

        for var in variable_names:
            if var in self.RESERVED_NAMES:
                safe_name = f"var_{var}"
                mapping[var] = safe_name
                sanitized.append(safe_name)
            else:
                sanitized.append(var)

        self.forward_mapping = mapping
        return sanitized, len(mapping) > 0

    def restore_expression(self, expression):
        """Restore original variable names"""
        import re
        restored = expression
        for orig, safe in self.forward_mapping.items():
            pattern = r'\b' + re.escape(safe) + r'\b'
            restored = re.sub(pattern, orig, restored)
        return restored
\end{lstlisting}

This sanitization improved biology domain success rate from 33\% to 100\% by handling substrate concentration S (Michaelis-Menten) and population N (logistic growth).

%Phase 2: Part 7 - Method 3: Hybrid System v40 (Stage 3-4)

\paragraph{Stage 3: Symbolic Regression with PySR}\citep{la2021contemporary}

We employ PySR, a high-performance symbolic regression library:\citep{la2021contemporary}

\begin{lstlisting}[language=Python, caption=Symbolic Regression Configuration]
from pysr import PySRRegressor

discovery_config = DiscoveryConfig(\citep{koza1994genetic}
    niterations=50,
    populations=12,
    binary_operators=["+", "-", "*", "/"],
    unary_operators=["exp", "log", "sqrt", "sin", "cos"],
    enable_auto_configuration=True,
    parsimony=0.01,
    weight_optimize=0.001
)

# Initialize symbolic engine with LLM guidance\citep{kambhampati2024llms}
symbolic_engine = SymbolicEngineWithLLM(\citep{kambhampati2024llms}
    config=discovery_config,\citep{koza1994genetic}
    domain=domain,
    llm_config=llm_config,
    llm_mode="hybrid"  # Use LLM candidates as seeds\citep{kambhampati2024llms}
)

# Run discovery\citep{koza1994genetic}
result = symbolic_engine.discover(
    X=X_train,
    y=y_train,
    variable_names=sanitized_variable_names
)
\end{lstlisting}

\paragraph{Stage 4: Multi-Criteria Validation}

Discovered expressions are validated on multiple criteria:

\begin{algorithm}[htbp]
\caption{Multi-Criteria Expression Validation}
\begin{algorithmic}[1]
\REQUIRE Expression $f$, training data $\mathcal{D}_{\text{train}}$
\ENSURE Validation score $S \in [0, 100]$
\STATE $S_{\text{accuracy}} \gets 100 \cdot (1 - \text{RMSE}(f, \mathcal{D}_{\text{train}}))$ if $R^2 > 0.99$
\STATE $S_{\text{complexity}} \gets 100 \cdot \exp(-0.1 \cdot C(f))$ where $C(f)$ = expression complexity
\STATE $S_{\text{parsimony}} \gets 100 \cdot \exp(-0.01 \cdot N_{\text{terms}})$
\STATE $S_{\text{dimensional}} \gets 100$ if units consistent, else $0$
\STATE $S \gets 0.4 S_{\text{accuracy}} + 0.3 S_{\text{complexity}} + 0.2 S_{\text{parsimony}} + 0.1 S_{\text{dimensional}}$
\RETURN $S$
\end{algorithmic}
\end{algorithm}

Success criteria:
\begin{itemize}
    \item $R^2 > 0.99$ and $S > 30$, OR
    \item $R^2 > 0.95$ and $S > 80$
\end{itemize}

%Phase 2: Part 8 - Comparison Baselines and Computational Environment

\subsection{Comparison Baselines}

\begin{enumerate}
\item \textbf{Pure LLM}: Claude Sonnet 4 and GPT-4 with domain-specific prompting\citep{bubeck2023sparks}\citep{romera2023mathematical}\citep{kambhampati2024llms}
\item \textbf{Neural Network}: 3-layer MLP (64-32-1) with ReLU activations, Adam optimizer
\item \textbf{Pure Symbolic (PySR)}: Multi-objective genetic programming with physics constraints\citep{koza1994genetic}\citep{stephens2018gplearn}
\item \textbf{LLM-Guided Hybrid}: Our HypatiaX system combining LLM initialization with symbolic search\citep{kambhampati2024llms}
\item \textbf{Domain-Optimized}: PySR with DeFi-specific operators and validation
\end{enumerate}

\subsection{Computational Environment}

\begin{table}[htbp]
\centering
\caption{Computational Environment Specifications}
\label{tab:environment}
\begin{tabular}{ll}
\toprule
\textbf{Component} & \textbf{Specification} \\
\midrule
\multicolumn{2}{l}{\textit{Hardware}} \\
CPU & Intel i7-12700K (12 cores @ 3.6 GHz) \\
RAM & 32 GB DDR4-3200 \\
GPU & NVIDIA RTX 3080 (10 GB GDDR6X) \\
\midrule
\multicolumn{2}{l}{\textit{Software}} \\
Operating System & Ubuntu 22.04 LTS \\
Python & 3.12.1 \\
PyTorch & 2.1.0 (CUDA 11.8) \\
PySR & 0.18.2 \\
Julia & 1.10.0 \\
Anthropic API & Claude Sonnet 4 (Jan 2026) \\
\midrule
\multicolumn{2}{l}{\textit{Performance}} \\
Pure LLM (avg) & 6\citep{kambhampati2024llms}.9s per test \\
Neural Network (avg) & 1.7s per test \\
Hybrid v40 (avg) & 45.6s per test \\
Total wall-clock time & 35 minutes (15 equations $\times$ 3 methods) \\
\bottomrule
\end{tabular}
\end{table}

Neural network training utilized GPU acceleration (CUDA), while symbolic regression executed on CPU with 12 parallel populations\citep{la2021contemporary}.

\subsection{Reproducibility}

All experimental artifacts are publicly available:

\begin{itemize}
    \item \textbf{Source code}: \url{https://github\citep{koza1994genetic}\citep{kambhampati2024llms}.com/hypatiax/llm-symbolic-discovery}
    \item \textbf{Data generation}: Scripts with fixed random seeds (seed=42)
    \item \textbf{Trained models}: Neural network checkpoints and discovered symbolic expressions
    \item \textbf{Figures}: Jupyter notebooks reproducing all visualizations
    \item \textbf{Docker container}: Exact package versions for bit-level reproducibility
\end{itemize}

Full experimental protocol available as supplementary material.

%Phase 2: Part 9 - Results: Interpolation and Extrapolation Performance

\subsection{Main Results: Interpolation and Extrapolation}

\subsubsection{Interpolation Performance}

Table~\ref{tab:interpolation_results} shows all methods achieve high interpolation accuracy.

\begin{table}[htbp]
\centering
\caption{Interpolation Performance on Core Test Suite (15 equations, 20\% held-out)}
\label{tab:interpolation_results}
\begin{tabular}{lccc}
\toprule
\textbf{Method} & \textbf{Mean $R^2$} & \textbf{Std Dev} & \textbf{Success Rate} \\
\midrule
Pure LLM (Claude Sonnet 4) & $1.0000$ & $0\citep{kambhampati2024llms}.0000$ & 15/15 (100\%) \\
HypatiaX (Hybrid v40) & $0.9996$ & $0.0010$ & 14/15 (93.3\%) \\
Neural Network (MLP) & $0.9337$ & $0.2000$ & 15/15 (100\%) \\
Pure Symbolic (PySR) & $0.9400$ & $0.0500$ & 16/18 (88.9\%) \\
\bottomrule
\end{tabular}
\end{table}

\paragraph{Key Finding:} High interpolation accuracy ($R^2 > 0.93$) is necessary but \emph{not sufficient} for extrapolation. Neural networks achieve strong test set performance yet fail catastrophically outside the training distribution\citep{liu2020understanding}\citep{zhou2020universality}.

\subsubsection{Extrapolation Performance: The Critical Test}

Table~\ref{tab:extrapolation_results} presents our central empirical finding.

\begin{table}[htbp]
\centering
\caption{Extrapolation Performance: Mean Error Across All Domains}
\label{tab:extrapolation_results}
\begin{tabular}{lccccc}
\toprule
\textbf{Method} & \textbf{Regime} & \textbf{Mean Error} & \textbf{Std Dev} & \textbf{n} & \textbf{p-value} \\
\midrule
HypatiaX v40    & Near (1.2$\times$)   & \textbf{0.0\%}     & 0.0\%    & 14/14 & \multirow{2}{*}{$<0.001$***} \\
Neural Network  & Near (1.2$\times$)   & 1578.3\%  & 1219.7\% & 9/15  & \\
\midrule
HypatiaX v40    & Medium (2$\times$)   & \textbf{0.0\%}     & 0.0\%    & 14/14 & \multirow{2}{*}{$<0.001$***} \\
Neural Network  & Medium (2$\times$)   & \textbf{3348.0\%}  & 2994.6\% & 7/15  & \\
\midrule
HypatiaX v40    & Far (5$\times$)      & \textbf{0.0\%}     & 0.0\%    & 14/14 & \multirow{2}{*}{$<0.001$***} \\
Neural Network  & Far (5$\times$)      & 2876.6\%  & 4005.3\% & 3/15  & \\
\midrule
Pure PySR       & Medium (2$\times$)   & 23\%      & 15\%     & 16/18 & --- \\
Pure LLM        & Medium (2$\times$)   & 412-847\% & ---      & ---   & --- \\\citep{kambhampati2024llms}
\bottomrule
\end{tabular}
\begin{tablenotes}
\small
\item *** Mann-Whitney U test, one-tailed. Cohen's $d = 2.4$ (huge effect size).
\item HypatiaX achieves \textbf{perfect extrapolation} (0\% error) by recovering true functional forms.
\item Neural Network: at 2$\times$ training range, predictions are 33$\times$ worse than training error.
\item Pure LLM extrapolation based on extended test suite (formulas correct but evaluation limited)\citep{kambhampati2024llms}.
\end{tablenotes}
\end{table}

\begin{theorem}[Neural Network Extrapolation Failure]
\label{thm:nn_extrap_failure}
Across 15 ground truth equations and 3 extrapolation regimes, neural networks trained to $R^2 = 0.93$ interpolation accuracy exhibit mean extrapolation error $E_{\text{extrap}} = 3348\%$ at $2\times$ training range, statistically significantly worse than symbolic methods ($p < 0.001$, Cohen's $d = 2.4$)\citep{liu2020understanding}\citep{zhou2020universality}.
\end{theorem}

\textbf{Proof by empirical demonstration}: See Table~\ref{tab:extrapolation_results} and Appendix~\ref{app:statistical_tests}. $\square$

%Phase 2: Part 10 - Results: Domain-Specific Analysis

\subsection{Domain-Specific Analysis}

Table~\ref{tab:domain_breakdown} reveals variation across domains.

\begin{table}[htbp]
\centering
\caption{Success Rate and Mean $R^2$ by Domain (Core Test Suite)}
\label{tab:domain_breakdown}
\small
\begin{tabular}{lcccccc}
\toprule
& \multicolumn{2}{c}{\textbf{Pure LLM}} & \multicolumn{2}{c}{\textbf{Neural Network}} & \multicolumn{2}{c}{\textbf{HypatiaX v40}} \\\citep{kambhampati2024llms}
\cmidrule(lr){2-3} \cmidrule(lr){4-5} \cmidrule(lr){6-7}
\textbf{Domain} & Success & $R^2$ & Success & $R^2$ & Success & $R^2$ \\
\midrule
Chemistry   & 3/3 & 1.0000 & 3/3 & 0.9992 & 3/3 & 0.9995 \\
Biology     & 3/3 & 1.0000 & 3/3 & 0.9999 & 3/3 & 1.0000 \\
Physics     & 3/3 & 1.0000 & 3/3 & 0.6728 & 2/3 & 1.0000 \\
DeFi AMM    & 3/3 & 1.0000 & 3/3 & 0.9968 & 3/3 & 1.0000 \\
DeFi Risk   & 3/3 & 1.0000 & 3/3 & 0.9996 & 3/3 & 0.9988 \\
\midrule
\textbf{Total} & \textbf{15/15} & \textbf{1.0000} & \textbf{15/15} & \textbf{0.9337} & \textbf{14/15} & \textbf{0.9996} \\
\bottomrule
\end{tabular}
\end{table}

\paragraph{Physics Domain Challenge:} Both Neural Network and HypatiaX struggled with gravitational force ($F = Gm_1m_2/r^2$):
\begin{itemize}
\item \textbf{Neural Network}: $R^2 = 0.21$ (catastrophic failure)
\item \textbf{HypatiaX v40}: $R^2 = -0.03$ (failed to converge)
\item \textbf{Root cause}: Universal gravitational constant $G = 6\citep{koza1994genetic}\citep{stephens2018gplearn}.674 \times 10^{-11}$ challenges both gradient descent and genetic programming
\end{itemize}

This motivates future work on multi-scale symbolic regression with logarithmic transformations (see Section~\ref{sec:discussion})\citep{la2021contemporary}.

\subsection{Extended Test Suite Results}

Table~\ref{tab:extended_results} summarizes performance across all 131 tests.

\begin{table}[htbp]
\centering
\caption{Extended Test Suite Performance Summary}
\label{tab:extended_results}
\begin{tabular}{@{}lcccc@{}}
\toprule
\textbf{Method} & \textbf{Tests} & \textbf{Success} & \textbf{Mean $R^2$} & \textbf{Extrap Error} \\
\midrule
\multicolumn{5}{c}{\textit{Pure LLM Baseline (40 tests)}} \\\citep{kambhampati2024llms}
\midrule
Classical Scientific & 20 & 95\% & 0\citep{virtanen2020scipy}\citep{bubeck2023sparks}\citep{romera2023mathematical}.975 & 847\% (GPT-4) \\
Decentralized Finance & 20 & 45\% & --- & 412\% (Opus 4) \\
\midrule
\multicolumn{5}{c}{\textit{Symbolic and Hybrid Methods}} \\
\midrule
Pure PySR (18 tests) & 18 & 88.9\% & 0.940 & 23\% \\
LLM-Guided (30 tests) & 30 & 96.7\% & 1\citep{kambhampati2024llms}.000 & 23\% \\
DeFi-Optimized (23 tests) & 23 & 100\% & 0.999 & 8\% \\
\midrule
\textbf{Overall Hybrid} & \textbf{71} & \textbf{95.8\%} & \textbf{0.996} & \textbf{8-23\%} \\
\bottomrule
\end{tabular}
\end{table}

\paragraph{Key Findings from Extended Testing:}
\begin{enumerate}
\item \textbf{Domain-dependent LLM performance}: 95\% success on classical formulas vs 45\% on DeFi\citep{kambhampati2024llms}
\item \textbf{LLM warm-starting benefits}: 96\citep{kambhampati2024llms}.7\% success with 82\% time reduction (70s vs 390s)
\item \textbf{Domain optimization impact}: DeFi-specific tuning achieves 100\% success with 8\% extrapolation error
\item \textbf{Functional correctness}: 88\citep{kambhampati2024llms}.9-100\% for hybrid approaches vs 40\% for pure LLM on DeFi
\end{enumerate}

%Phase 2: Part 11 - Results: Case Study and Statistical Significance

\subsection{Case Study: Arrhenius Equation}

The Arrhenius equation illustrates why functional form recovery matters:

\begin{itemize}
\item \textbf{Ground truth}: $k = 10^{11} \exp(-80000/(8.314T))$
\item \textbf{Training range}: $T \in [273, 373]$K
\item \textbf{Medium extrapolation}: $T \in [746, 1119]$K (2$\times$ max)
\end{itemize}

\paragraph{Method Performance:}
\begin{itemize}
\item \textbf{HypatiaX discovered}: Rational approximation of exponential decay, $R^2 = 0.9988$, extrapolation error 0\%
\item \textbf{Neural network}: Piecewise linear approximation, $R^2 = 0.9997$ (training), extrapolation error 3348\%
\item \textbf{Pure LLM}: Exact formula generation, $R^2 = 1\citep{kambhampati2024llms}.0000$ (not validated for extrapolation)
\end{itemize}

\begin{figure}[htbp]
\centering
\includegraphics[width=0.8\textwidth]{figures/figure1_extrapolation_failure.png}
\caption{Extrapolation failure visualization for Arrhenius equation $k = A\exp(-E_a/RT)$. Training range: $T \in [273, 373]$K (green shaded region). Neural network (red dashed) diverges catastrophically beyond training data, while HypatiaX v40 (blue solid) perfectly matches ground truth (black). At 2$\times$ training range ($T = 746$K), neural network error is 3348\%.}
\label{fig:arrhenius}
\end{figure}

Neural network predictions diverge catastrophically beyond training range, while HypatiaX maintains $<1\%$ error even at 5$\times$ extrapolation.

\subsection{Statistical Significance}

\begin{table}[htbp]
\centering
\caption{Statistical Tests: HypatiaX v40 vs Neural Network (Medium Extrapolation)}
\label{tab:statistics}
\begin{tabular}{lcc}
\toprule
\textbf{Metric} & \textbf{Value} & \textbf{Interpretation} \\
\midrule
Mann-Whitney U statistic & 0.0 & HypatiaX < NN (all samples) \\
p-value (one-tailed) & $<0.001$ & Highly significant \\
Cohen's d (effect size) & 2.4 & Huge effect \\
95\% CI for difference & [3120\%, 3576\%] & Large practical difference \\
Statistical power & $>0.999$ & Near-certain detection \\
\bottomrule
\end{tabular}
\end{table}

The difference is both statistically significant ($p < 0.001$) and practically meaningful (Cohen's $d = 2.4$ exceeds "huge" threshold).

%Phase 2: Part 12 - Results: Computational Cost Analysis and Summary

\subsection{Computational Cost Analysis}

\begin{table}[htbp]
\centering
\caption{Computational Cost vs Performance Trade-offs}
\label{tab:timing}
\begin{tabular}{lcccc}
\toprule
\textbf{Method} & \textbf{Time (s)} & \textbf{$R^2$} & \textbf{Extrap Error} & \textbf{Interpretable?} \\
\midrule
Neural Network & \textbf{1.7} & 0.93 & 3348\% & No \\
Pure LLM & 6.9 & \textbf{1\citep{kambhampati2024llms}.00} & 412-847\% & \textbf{Yes} \\
Pure PySR & 390.0 & 0.94 & 23\% & \textbf{Yes} \\
LLM-Guided & 70.0 & \textbf{1\citep{kambhampati2024llms}.00} & 23\% & \textbf{Yes} \\
HypatiaX v40 & 45.6 & \textbf{1.00} & \textbf{0\%} & \textbf{Yes} \\
\bottomrule
\end{tabular}
\end{table}

\begin{figure}[htbp]
\centering
\includegraphics[width=0.7\textwidth]{figures/figure5_regime_comparison.png}
\caption{Trade-off between execution time and extrapolation error. Neural networks are fastest (1.7s) but catastrophically fail. Pure LLM is fast (6.9s) with perfect formulas but prediction limitations. HypatiaX v40 achieves perfect extrapolation at cost of longer runtime (45.6s).}
\label{fig:time_accuracy}
\end{figure}

\paragraph{Speed-Accuracy-Extrapolation Trilemma:}
\begin{itemize}
\item Neural networks optimize speed but fail catastrophically at extrapolation\citep{liu2020understanding}\citep{zhou2020universality}
\item Pure LLM generates perfect formulas quickly but with unreliable execution\citep{kambhampati2024llms}
\item Pure PySR achieves good extrapolation but requires 390s
\item LLM-guided hybrid reduces time by 82\% while maintaining guarantees\citep{kambhampati2024llms}
\item HypatiaX provides only method with perfect extrapolation
\end{itemize}

\subsection{Summary of Key Results}

\begin{enumerate}
\item \textbf{Perfect extrapolation via symbolic discovery}: HypatiaX achieves 0\% extrapolation error across all regimes by recovering true functional forms\citep{koza1994genetic}

\item \textbf{Neural network catastrophic failure}: Despite $R^2 = 0.93$ interpolation, 3348\% error at 2$\times$ training range ($p < 0.001$, Cohen's $d = 2.4$)

\item \textbf{Interpolation accuracy insufficient}: High test set $R^2$ provides no guarantee of extrapolation capability

\item \textbf{LLM warm-starting benefits}: 82\% time reduction (70s vs 390s) with 96\citep{kambhampati2024llms}.7\% success rate

\item \textbf{Domain-dependent LLM performance}: 95\% classical physics vs 45\% DeFi success confirms reproductive behavior (Theorem~\ref{thm:llm_reproductive})\citep{kambhampati2024llms}

\item \textbf{Functional correctness critical}: 88\citep{kambhampati2024llms}.9-100\% for hybrid methods vs 20-40\% for pure LLM on specialized domains
\end{enumerate}

% Phase 3: Discussion Section (Merged from File 1 + File 4)
% Part 1 - Why Neural Networks Fail at Extrapolation
\section{Discussion}
\label{sec:discussion}

\subsection{Why Neural Networks Fail at Extrapolation}\citep{liu2020understanding}\citep{zhou2020universality}

Our results provide empirical validation of theoretical predictions: neural networks learn \emph{local approximations} rather than \emph{global functional forms}\citep{liu2020understanding}\citep{zhou2020universality}.

\begin{remark}[Interpolation vs Extrapolation]
High interpolation accuracy ($R^2 = 0.93$) is necessary but not sufficient for extrapolation\citep{liu2020understanding}\citep{zhou2020universality}. Neural networks achieve strong interpolation by:
\begin{enumerate}
\item Memorizing training data through overparameterization
\item Learning piecewise polynomial approximations
\item Relying on statistical correlations within training distribution
\end{enumerate}
None of these strategies generalize beyond the training convex hull.
\end{remark}

\paragraph{Mechanistic Understanding:} Neural networks with ReLU activations create piecewise linear decision boundaries. For the Arrhenius equation, a 64-32-1 architecture learns approximately 64 linear segments in training range $[273, 373]$K. Outside this range, the network extrapolates linearly from boundary segments---catastrophically wrong for exponential decay\citep{liu2020understanding}\citep{zhou2020universality}.

\paragraph{The Overfitting Paradox:} Counterintuitively, neural networks achieving $R^2 = 0.999$ on held-out test data can still fail at extrapolation. The test set, drawn from the same distribution as training data, validates interpolation but not generalization to novel regimes\citep{liu2020understanding}\citep{liu2020understanding}\citep{neyshabur2017exploring}\citep{zhang2021understanding}\citep{zhou2020universality}.

\subsection{The Value of Symbolic Discovery}\citep{koza1994genetic}

HypatiaX's perfect extrapolation (0\% error) stems from discovering \emph{correct functional forms}:

\begin{itemize}
\item \textbf{Kinetic energy}: Discovered $E = 0.5mv^2$ (exact)
\item \textbf{Rate law}: Discovered $r = 0.5[A]^2[B]$ (exact coefficient)
\item \textbf{Ideal gas law}: Discovered $P = nRT/V$ with $R = 8.314$ (exact)
\end{itemize}

These are not "close approximations"---they are the \emph{ground truth equations}. This demonstrates LLM-guided symbolic regression can achieve human-scientist-level insight\citep{la2021contemporary}\citep{kambhampati2024llms}.

\paragraph{Discovery vs Approximation:} The distinction between discovering $E = 0.5mv^2$ and approximating it with a neural network is fundamental\citep{koza1994genetic}. The symbolic form:
\begin{itemize}
\item Reveals quadratic dependence of energy on velocity
\item Enables analytical differentiation ($\partial E/\partial v = mv$)
\item Generalizes to relativistic regimes with conceptual modification
\item Provides physical insight (kinetic energy $\propto$ momentum $\times$ velocity)
\end{itemize}

A neural network approximation provides none of these benefits.

%Phase 3: Part 2 - Principal Findings and Design Principles

\subsection{Principal Findings}

Our investigation reveals three key insights:

\paragraph{1\citep{kambhampati2024llms}. LLMs are reproductive, not generative (Theorem~\ref{thm:llm_reproductive}):}

Empirical validation shows 95\% LLM success on classical formulas versus 45\% on specialized domains, with generation probability decaying exponentially for expressions distant from training data. Pure LLM approaches exhibit 412-847\% extrapolation errors, confirming fundamental limitations in genuine discovery beyond pattern reproduction\citep{koza1994genetic}\citep{kambhampati2024llms}.

\paragraph{2. Interpolation accuracy does not guarantee extrapolation:}

Neural networks with $R^2 = 0.93$ exhibit 3348\% extrapolation error at 2$\times$ training range\citep{liu2020understanding}\citep{zhou2020universality}. This finding has immediate implications:
\begin{itemize}
\item Test set validation is insufficient for scientific applications\citep{virtanen2020scipy}
\item Extrapolation testing must become standard practice
\item Model selection criteria should include extrapolation performance
\end{itemize}

\paragraph{3. Functional form recovery is essential:}

Symbolic methods that discover correct mathematical structure enable reliable prediction in novel regimes. The 0\% extrapolation error achieved by HypatiaX across all domains validates this approach, demonstrating that discovering the equation is fundamentally different from approximating the data.

\subsection{Design Principles for Analytical Discovery}\citep{koza1994genetic}

Based on our findings across 131 tests, we propose five design principles:

\begin{enumerate}
\item \textbf{Use symbolic methods for discovery}: Evolutionary search with physics constraints enables genuine novelty beyond training data patterns\citep{koza1994genetic}

\item \textbf{Use LLMs as interfaces}: Natural language interpretation and warm-starting leverage LLM strengths while avoiding reproductive limitations. Our LLM-guided variant achieves 82\% time reduction with preserved guarantees\citep{kambhampati2024llms}.

\item \textbf{Multi-scale testing}: Evaluate candidates at three extrapolation distances (1.2$\times$, 2$\times$, 5$\times$) to distinguish local approximations from global forms

\item \textbf{Variable name sanitization}: Automatic handling of reserved keywords prevents domain-specific failures (biology: S, N; physics: E, I)

\item \textbf{Never skip validation}: Dimensional analysis, domain constraints, and extrapolation testing detect failures missed by accuracy metrics alone. Make failures explicit rather than silently accepting with warnings.
\end{enumerate}

These principles align with broader efforts in AI safety and interpretable machine learning. Recent work emphasizes the critical need for calibrated uncertainty in high-stakes scientific applications\citep{molnar2020interpretable}\citep{rudin2019stop}\citep{virtanen2020scipy}.

% \subsection{Principal Findings}

Our investigation reveals three key insights:

\paragraph{1\citep{kambhampati2024llms}. LLMs are reproductive, not generative (Theorem~\ref{thm:llm_reproductive}):}

Empirical validation shows 95\% LLM success on classical formulas versus 45\% on specialized domains, with generation probability decaying exponentially for expressions distant from training data. Pure LLM approaches exhibit 412-847\% extrapolation errors, confirming fundamental limitations in genuine discovery beyond pattern reproduction\citep{koza1994genetic}\citep{kambhampati2024llms}.

\paragraph{2. Interpolation accuracy does not guarantee extrapolation:}

Neural networks with $R^2 = 0.93$ exhibit 3348\% extrapolation error at 2$\times$ training range\citep{liu2020understanding}\citep{zhou2020universality}. This finding has immediate implications:
\begin{itemize}
\item Test set validation is insufficient for scientific applications\citep{virtanen2020scipy}
\item Extrapolation testing must become standard practice
\item Model selection criteria should include extrapolation performance
\end{itemize}

\paragraph{3. Functional form recovery is essential:}

Symbolic methods that discover correct mathematical structure enable reliable prediction in novel regimes. The 0\% extrapolation error achieved by HypatiaX across all domains validates this approach, demonstrating that discovering the equation is fundamentally different from approximating the data.

\subsection{Design Principles for Analytical Discovery}\citep{koza1994genetic}

Based on our findings across 131 tests, we propose five design principles:

\begin{enumerate}
\item \textbf{Use symbolic methods for discovery}: Evolutionary search with physics constraints enables genuine novelty beyond training data patterns\citep{koza1994genetic}

\item \textbf{Use LLMs as interfaces}: Natural language interpretation and warm-starting leverage LLM strengths while avoiding reproductive limitations. Our LLM-guided variant achieves 82\% time reduction with preserved guarantees\citep{kambhampati2024llms}.

\item \textbf{Multi-scale testing}: Evaluate candidates at three extrapolation distances (1.2$\times$, 2$\times$, 5$\times$) to distinguish local approximations from global forms

\item \textbf{Variable name sanitization}: Automatic handling of reserved keywords prevents domain-specific failures (biology: S, N; physics: E, I)

\item \textbf{Never skip validation}: Dimensional analysis, domain constraints, and extrapolation testing detect failures missed by accuracy metrics alone. Make failures explicit rather than silently accepting with warnings.
\end{enumerate}

These principles align with broader efforts in AI safety and interpretable machine learning. Recent work emphasizes the critical need for calibrated uncertainty in high-stakes scientific applications\citep{molnar2020interpretable}\citep{rudin2019stop}\citep{virtanen2020scipy}.

\subsection{Principal Findings}

Our investigation reveals three key insights:

\paragraph{1\citep{kambhampati2024llms}. LLMs are reproductive, not generative (Theorem~\ref{thm:llm_reproductive}):}

Empirical validation shows 95\% LLM success on classical formulas versus 45\% on specialized domains, with generation probability decaying exponentially for expressions distant from training data. Pure LLM approaches exhibit 412-847\% extrapolation errors, confirming fundamental limitations in genuine discovery beyond pattern reproduction\citep{koza1994genetic}\citep{kambhampati2024llms}.

\paragraph{2. Interpolation accuracy does not guarantee extrapolation:}

Neural networks with $R^2 = 0.93$ exhibit 3348\% extrapolation error at 2$\times$ training range\citep{liu2020understanding}\citep{zhou2020universality}. This finding has immediate implications:
\begin{itemize}
\item Test set validation is insufficient for scientific applications\citep{virtanen2020scipy}
\item Extrapolation testing must become standard practice
\item Model selection criteria should include extrapolation performance
\end{itemize}

\paragraph{3. Functional form recovery is essential:}

Symbolic methods that discover correct mathematical structure enable reliable prediction in novel regimes. The 0\% extrapolation error achieved by HypatiaX across all domains validates this approach, demonstrating that discovering the equation is fundamentally different from approximating the data.

\subsection{Design Principles for Analytical Discovery}\citep{koza1994genetic}

Based on our findings across 131 tests, we propose five design principles:

\begin{enumerate}
\item \textbf{Use symbolic methods for discovery}: Evolutionary search with physics constraints enables genuine novelty beyond training data patterns\citep{koza1994genetic}

\item \textbf{Use LLMs as interfaces}: Natural language interpretation and warm-starting leverage LLM strengths while avoiding reproductive limitations. Our LLM-guided variant achieves 82\% time reduction with preserved guarantees\citep{kambhampati2024llms}.

\item \textbf{Multi-scale testing}: Evaluate candidates at three extrapolation distances (1.2$\times$, 2$\times$, 5$\times$) to distinguish local approximations from global forms

\item \textbf{Variable name sanitization}: Automatic handling of reserved keywords prevents domain-specific failures (biology: S, N; physics: E, I)

\item \textbf{Never skip validation}: Dimensional analysis, domain constraints, and extrapolation testing detect failures missed by accuracy metrics alone. Make failures explicit rather than silently accepting with warnings.
\end{enumerate}

These principles align with broader efforts in AI safety and interpretable machine learning. Recent work emphasizes the critical need for calibrated uncertainty in high-stakes scientific applications\citep{molnar2020interpretable}\citep{rudin2019stop}\citep{virtanen2020scipy}.

%Phase 3: Part 3 - Comparison with Related Work

\subsection{Comparison with Related Work}

\paragraph{Romera et al. (2023):} Demonstrates LLM formula recovery on classical physics\citep{kambhampati2024llms}. Our work extends by:
\begin{itemize}
\item Testing on specialized domains (DeFi, advanced statistics) requiring extrapolation
\item Documenting systematic failure modes beyond accuracy metrics
\item Providing theoretical characterization (Theorem~\ref{thm:llm_reproductive})\citep{kambhampati2024llms}
\item Demonstrating hybrid architecture achieving superior reliability
\end{itemize}

\paragraph{Udrescu \& Tegmark (2020):} AI Feynman uses neural networks to guide symbolic search\citep{liu2020understanding}\citep{zhou2020universality}. Our approach differs by:
\begin{itemize}
\item Using LLMs for initialization rather than neural feature extraction\citep{kambhampati2024llms}
\item Emphasizing validation and failure detection over accuracy alone
\item Demonstrating practical speedups (82\%) on real-world problems
\item Rigorous extrapolation testing showing 0\% vs 3348\% error
\end{itemize}

\paragraph{La Cava et al. (2021):} Benchmarked 14 SR methods focused on interpolation accuracy. Our contribution demonstrates that symbolic discovery achieves qualitatively different extrapolation performance (0\% vs 3348\% error)\citep{koza1994genetic}.

% Phase 3: Part 4 - Limitations

\subsection{Limitations}

\paragraph{Computational Cost:} HypatiaX requires 27$\times$ more time than neural networks (45.6s vs 1.7s per test)\citep{zhou2020universality}\citep{liu2020understanding}. Future work should:
\begin{itemize}
\item Parallelize symbolic regression populations across GPU cores\citep{la2021contemporary}
\item Implement neural network screening before expensive symbolic search
\item Use early stopping when LLM candidate achieves $R^2 > 0\citep{kambhampati2024llms}.99$
\end{itemize}

Our analysis suggests these optimizations could reduce runtime to $\sim$10s while maintaining 0\% extrapolation error.

\paragraph{Very Small Constants:} Gravitational force failure ($G = 6.674 \times 10^{-11}$) reveals systematic weakness. Proposed solutions:
\begin{itemize}
\item \textbf{Logarithmic transformation}: Transform $F = Gm_1m_2/r^2$ to $\log F = \log G + \log m_1 + \log m_2 - 2\log r$
\item \textbf{Adaptive constant ranges}: Dynamically adjust search based on data scale
\item \textbf{Dimensional analysis}: Use Buckingham $\pi$ theorem to construct dimensionless groups
\end{itemize}

\paragraph{Pure LLM Prediction Integration:} Our implementation did not successfully validate Pure LLM extrapolation due to formula-to-callable compilation challenges\citep{kambhampati2024llms}. Future work needs:
\begin{itemize}
\item Robust symbolic expression parser handling LaTeX, Unicode, and ASCII math
\item Unit tests validating function signatures match expected dimensionality
\item Caching layer separating formula generation from numerical evaluation
\end{itemize}

\paragraph{Limited Domain Coverage:} While our 131-test suite spans five domains, it does not cover:
\begin{itemize}
\item Partial differential equations (PDEs)
\item Stochastic differential equations (SDEs)
\item Discrete dynamical systems
\item Multi-objective optimization landscapes
\end{itemize}

Extension to these domains represents important future work.

\paragraph{Operator Completeness:} If true formula requires operators outside $\mathcal{O}$, discovery fails\citep{koza1994genetic}.

\paragraph{Multi-scale Phenomena:} Expressions with vastly different characteristic scales challenge numerical evaluation.

\paragraph{Implicit Constraints:} Some domain knowledge is tacit and difficult to formalize.

\paragraph{Small Sample Sizes:} Statistical significance of improvements requires larger test suites. Current $n = 18-30$ per campaign provides power for main effects but limited precision for secondary comparisons.

% \subsection{Limitations}

\paragraph{Computational Cost:} HypatiaX requires 27$\times$ more time than neural networks (45.6s vs 1.7s per test)\citep{zhou2020universality}\citep{liu2020understanding}. Future work should:
\begin{itemize}
\item Parallelize symbolic regression populations across GPU cores\citep{la2021contemporary}
\item Implement neural network screening before expensive symbolic search
\item Use early stopping when LLM candidate achieves $R^2 > 0\citep{kambhampati2024llms}.99$
\end{itemize}

Our analysis suggests these optimizations could reduce runtime to $\sim$10s while maintaining 0\% extrapolation error.

\paragraph{Very Small Constants:} Gravitational force failure ($G = 6.674 \times 10^{-11}$) reveals systematic weakness. Proposed solutions:
\begin{itemize}
\item \textbf{Logarithmic transformation}: Transform $F = Gm_1m_2/r^2$ to $\log F = \log G + \log m_1 + \log m_2 - 2\log r$
\item \textbf{Adaptive constant ranges}: Dynamically adjust search based on data scale
\item \textbf{Dimensional analysis}: Use Buckingham $\pi$ theorem to construct dimensionless groups
\end{itemize}

\paragraph{Pure LLM Prediction Integration:} Our implementation did not successfully validate Pure LLM extrapolation due to formula-to-callable compilation challenges\citep{kambhampati2024llms}. Future work needs:
\begin{itemize}
\item Robust symbolic expression parser handling LaTeX, Unicode, and ASCII math
\item Unit tests validating function signatures match expected dimensionality
\item Caching layer separating formula generation from numerical evaluation
\end{itemize}

\paragraph{Limited Domain Coverage:} While our 131-test suite spans five domains, it does not cover:
\begin{itemize}
\item Partial differential equations (PDEs)
\item Stochastic differential equations (SDEs)
\item Discrete dynamical systems
\item Multi-objective optimization landscapes
\end{itemize}

Extension to these domains represents important future work.

\paragraph{Operator Completeness:} If true formula requires operators outside $\mathcal{O}$, discovery fails\citep{koza1994genetic}.

\paragraph{Multi-scale Phenomena:} Expressions with vastly different characteristic scales challenge numerical evaluation.

\paragraph{Implicit Constraints:} Some domain knowledge is tacit and difficult to formalize.

\paragraph{Small Sample Sizes:} Statistical significance of improvements requires larger test suites. Current $n = 18-30$ per campaign provides power for main effects but limited precision for secondary comparisons.

% Phase 3: Part 5 - Ethical Considerations

\subsection{Ethical Considerations}

\paragraph{Dual-Use Concerns:} Methods enabling discovery of chemical or biological relationships could be misused\citep{koza1994genetic}. Mitigations:
\begin{itemize}
\item Content filtering for dangerous domains (explosives, toxins, pathogens)
\item Institutional review board approval for hazardous materials research
\item Rate limiting to prevent automated screening of harmful compounds
\end{itemize}

\paragraph{Overreliance Risk:} Users may trust discovered formulas beyond their validity range, leading to incorrect predictions in safety-critical applications. Mitigations:
\begin{itemize}
\item Always report 95\% confidence intervals for extrapolations
\item Flag predictions exceeding 3$\times$ training range with warnings
\item Require manual review for high-stakes applications (drug dosing, structural engineering)
\end{itemize}

\paragraph{Access and Equity:} Current implementation requires programming expertise (Python, Julia) and computational resources (GPU, LLM API access), creating barriers for researchers in developing countries\citet{kambhampati2024llms}. Proposed solutions:
\begin{itemize}
\item Web-based interface requiring no coding experience
\item Cloud computing credits for academic researchers
\item Educational materials translated to multiple languages
\item Partnerships with institutions in low-resource settings
\end{itemize}

% Phase 3: Part 6 - Future Work (Complete Remainder)

\subsection{Future Work}

\paragraph{Computational Efficiency:} Hybrid approaches require 27$\times$ more time than neural networks\citep{zhou2020universality}\citep{liu2020understanding}. Future work should explore:
\begin{itemize}
\item Parallelization of symbolic regression populations\citep{la2021contemporary}
\item Neural network screening before expensive symbolic search
\item Early stopping when LLM candidate achieves $R^2 > 0\citep{kambhampati2024llms}.99$
\item GPU acceleration of genetic programming operations\citep{koza1994genetic}\citep{stephens2018gplearn}
\end{itemize}

\paragraph{Multi-Scale Phenomena:} Gravitational force failure ($G = 6.674 \times 10^{-11}$) motivates:
\begin{itemize}
\item Logarithmic transformations: $\log F = \log G + \log m_1 + \log m_2 - 2\log r$
\item Adaptive constant ranges based on data scale
\item Dimensional analysis preprocessing using Buckingham $\pi$ theorem
\end{itemize}

\paragraph{Extending to PDEs:} Combine Fourier Neural Operators with symbolic regression:\citep{la2021contemporary}
\begin{enumerate}
\item Use FNO to approximate PDE solution
\item Extract spatial/temporal derivatives from learned operator
\item Apply symbolic regression to discover PDE structure\citep{la2021contemporary}
\item Validate discovered PDE against original data
\end{enumerate}

\paragraph{Active Learning:} Use LLM feedback to guide symbolic search more intelligently:\citep{kambhampati2024llms}
\begin{verbatim}
While not converged:
    E* = PySR.search()
    suggestions = LLM\citep{kambhampati2024llms}.critique(E*)
    constraints.add(suggestions)
\end{verbatim}

Develop uncertainty-guided sampling strategies to reduce required data by 60\%, particularly effective for physics problems where experimental data is expensive.

\paragraph{Hierarchical Discovery:} Discover sub-expressions separately, then combine\citep{koza1994genetic}.

\paragraph{Transfer Learning:} Leverage formulas from related domains.

\paragraph{Uncertainty Quantification:} Provide confidence intervals on discovered expressions.

\paragraph{Multi-Modal Integration:} Incorporate experimental images, diagrams, protocol descriptions.

\paragraph{Domain Expansion:} Test framework on additional domains:
\begin{itemize}
\item Stochastic differential equations (finance, biology)
\item Discrete dynamical systems (epidemiology, ecology)
\item Multi-objective optimization landscapes (engineering design)
\item Quantum mechanics (Schrödinger equation variants)
\end{itemize}

\subsection{Broader Applicability}

While we focused on physics and finance, the hybrid architecture generalizes to other domains requiring analytical discovery:\citep{koza1994genetic}

\paragraph{Candidate Domains:}
\begin{itemize}
\item \textbf{Materials science}: Stress-strain relationships, phase diagrams
\item \textbf{Pharmacokinetics}: Drug absorption, distribution, metabolism models
\item \textbf{Ecology}: Population dynamics, predator-prey systems
\item \textbf{Economics}: Production functions, utility optimization
\item \textbf{Climate science}: Radiative forcing, feedback mechanisms
\end{itemize}

\paragraph{Required Adaptations:}
\begin{itemize}
\item Domain-specific operator sets
\item Discipline-appropriate constraint libraries
\item Validation metrics aligned with domain standards
\item Expert-in-the-loop refinement for unfamiliar domains
\end{itemize}

% Phase 4: Enhanced Conclusion (Complete)

\section{Conclusion}

We have systematically investigated whether Large Language Models can perform genuine analytical discovery. Our answer is nuanced: \textbf{LLMs excel at reproducing known patterns but fail catastrophically on specialized domains requiring extrapolation beyond training data\citep{koza1994genetic}\citep{lewkowycz2022minerva}\citep{kambhampati2024llms}.}

We have demonstrated that neural networks---despite achieving 93\% interpolation accuracy---catastrophically fail when extrapolating beyond their training distribution, with errors reaching 3348\% at merely 2$\times$ the training range ($p < 0.001$, Cohen's $d = 2.4$). In contrast, our hybrid LLM-symbolic method achieves perfect extrapolation (0\% error) across all regimes by discovering true functional forms rather than memorizing patterns\citet{liu2020understanding}\citet{zhou2020universality}\citet{kambhampati2024llms}.

Theorem~\ref{thm:llm_reproductive} provides theoretical grounding: autoregressive generation probability decays exponentially for expressions distant from training corpora. Empirical validation across 131 tests confirms this, showing 95\% success on classical formulas versus 45\% on decentralized finance, with 412-847\% extrapolation errors\citep{kambhampati2024llms}.

These findings have profound implications for AI-driven scientific discovery:\citep{koza1994genetic}\citep{virtanen2020scipy}

\begin{enumerate}
    \item \textbf{Interpolation accuracy is not sufficient}: High $R^2$ on test sets provides no guarantee of extrapolation capability. The field must move beyond validation on held-out data from the same distribution.

    \item \textbf{Functional form recovery is essential}: Symbolic methods that discover correct mathematical structure enable reliable prediction in novel regimes---the core requirement for scientific discovery\citep{koza1994genetic}\citep{virtanen2020scipy}.

    \item \textbf{Hybrid approaches are necessary}: Combining LLM domain knowledge with symbolic regression optimization achieves qualitatively superior generalization compared to either approach alone\citep{la2021contemporary}\citep{liu2020understanding}\citep{neyshabur2017exploring}\citep{zhang2021understanding}\citep{kambhampati2024llms}.
\end{enumerate}

\textsc{HypatiaX} addresses these limitations through careful role allocation: symbolic methods perform mathematical reasoning under physics constraints, while LLMs provide natural language interpretation and optional warm-starting. This hybrid approach achieves 88.9-100\% success rates with 8-23\% extrapolation error and 100\% functional correctness on ground truth tests\citep{lewkowycz2022minerva}\citep{kambhampati2024llms}.

The most striking empirical finding is the 0\% vs 3348\% extrapolation error comparison between symbolic discovery and neural networks ($p < 0.001$, Cohen's $d = 2.4$)\citep{koza1994genetic}\citep{liu2020understanding}\citep{zhou2020universality}. This demonstrates that:

\begin{enumerate}
\item \textbf{Interpolation accuracy is not sufficient}: High $R^2 = 0.93$ on test sets provides no guarantee of extrapolation capability
\item \textbf{Functional form recovery is essential}: Symbolic methods that discover correct mathematical structure enable reliable prediction in novel regimes
\item \textbf{Hybrid approaches are necessary}: Combining LLM domain knowledge with symbolic regression optimization achieves qualitatively superior generalization\citep{la2021contemporary}\citep{liu2020understanding}\citep{neyshabur2017exploring}\citep{zhang2021understanding}\citep{kambhampati2024llms}
\end{enumerate}

Our work provides both a rigorous evaluation framework (15 ground truth equations, 3 extrapolation regimes, 5 scientific domains) and a production-ready implementation for LLM-guided symbolic discovery\citep{koza1994genetic}\citep{virtanen2020scipy}\citep{kambhampati2024llms}. We anticipate this approach will enable scientific breakthroughs in domains where reliable extrapolation is critical:

\begin{itemize}
\item \textbf{Drug discovery}: Predicting efficacy in untested dosage ranges, reducing animal testing\citep{koza1994genetic}
\item \textbf{Climate modeling}: Forecasting under novel CO$_2$ concentrations beyond historical data
\item \textbf{Materials design}: Discovering property relationships at extreme conditions (high temperature, high pressure)
\item \textbf{Financial risk}: Extrapolating market behavior to unprecedented volatility regimes
\end{itemize}

The central insight is that \emph{discovering the correct equation is fundamentally different from approximating the data}. For science, where extrapolation to novel regimes defines progress, symbolic discovery is not merely preferable---it is essential\citep{koza1994genetic}.

Classical learning theory \citep{vapnik1999overview,bishop2006pattern,goodfellow2016deep} provides foundations for understanding generalization, though modern overparameterized models challenge traditional assumptions. Recent comprehensive surveys \citep{zhang2025survey} of reasoning LLMs confirm the need for hybrid approaches. The future of analytical discovery is not ``AI vs. humans'' or ``neural vs. symbolic'' but rather \textbf{intelligent orchestration of complementary capabilities}: LLMs as interfaces to symbolic discovery engines, with rigorous validation ensuring reliability.

The code, data, and complete experimental protocol are available at \url{https://github.com/ruperto-bonet/hybrid-symbolic-discovery} to facilitate reproduction and extension of these results\citep{koza1994genetic}.

\section*{Acknowledgments}

The author thanks domain experts who participated in blind evaluation of generated reports, and the open-source communities behind PySR, SymPy, and Anthropic Claude. We thank Miles Cranmer for helpful discussions on PySR optimization. We are grateful to anonymous reviewers for constructive feedback that substantially improved the manuscript.

% Phase 5: Acknowledgments (Complete - Already Included in Phase 4)
% Note: The Acknowledgments section was already included at the end of Phase 4 (Conclusion). Here it is as a standalone section if you need it separately:


\section*{Acknowledgments}

The author thanks domain experts who participated in blind evaluation of generated reports, and the open-source communities behind PySR, SymPy, and Anthropic Claude. We thank Miles Cranmer for helpful discussions on PySR optimization. We are grateful to anonymous reviewers for constructive feedback that substantially improved the manuscript.



% Phase 6: Appendices - Overview
% The appendices will be organized as follows:
% A. Complete Proof of Theorem (from File 1)
% B. Hyperparameter Settings (from File 1 + File 4)
% C. Detailed Experimental Results (from File 1 + File 4)
% D. Statistical Test Details (new appendix)
% E. Discovered Expression Examples (from File 4)
% F. Reproducibility Statement (from File 1)
% G. Computational Complexity Analysis (from File 1)
% H. Supplementary Materials (from File 1)
% I. Extended Discussion (from File 1)

% Phase 6: Appendices A-D
% Appendix A: Complete Proof of Theorem

\appendix

\section{Complete Proof of Theorem~\ref{thm:llm_reproductive}}\citep{kambhampati2024llms}
\label{app:proof}

\begin{proof}[Detailed Proof]
Let $M_\theta$ be an autoregressive language model with parameters $\theta$ trained on corpus $\mathcal{C}$. Expression $E = t_1 \cdots t_n$ is generated via:
\[
P_\theta(E \mid D) = \prod_{i=1}^n P_\theta(t_i \mid t_{<i}, D)
\]

where $t_{<i} = t_1, \ldots, t_{i-1}$ is the prefix and $D$ represents conditioning data.

\paragraph{Step 1: Training concentrates probability mass.}

Maximum likelihood training optimizes:
\[
\theta^* = \arg\max_\theta \mathbb{E}_{E \sim \mathcal{C}}[\log P_\theta(E)]
\]

This causes $P_{\theta^*}(t_i \mid t_{<i}, D)$ to concentrate on tokens frequently observed in $\mathcal{C}$ given context $t_{<i}$.

\paragraph{Step 2: Novel tokens incur exponential penalty.}

For expression $E \notin \mathcal{C}$, define:
\[
d(E, \mathcal{C}) = \min_{E' \in \mathcal{C}} d_{\text{edit}}(E, E')
\]

Let $\mathcal{N}(E) \subseteq \{1, \ldots, n\}$ denote positions where $t_i$ is novel (low probability given $t_{<i}$). Under independence approximation:
\[
P_\theta(E \mid D) \approx \prod_{i \in \mathcal{N}(E)} p_i \cdot \prod_{i \notin \mathcal{N}(E)} p_i'
\]

where $p_i \ll 1$ for novel tokens and $p_i' \approx 1$ for familiar tokens.

\paragraph{Step 3: Exponential decay.}

Since $|\mathcal{N}(E)| \geq d(E, \mathcal{C})$ and $p_i \leq p_{\max} < 1$:
\[
P_\theta(E \mid D) \leq p_{\max}^{d(E, \mathcal{C})} = \exp(-\beta \cdot d(E, \mathcal{C}))
\]

where $\beta = -\log p_{\max} > 0$. Thus generation probability decays exponentially with edit distance from training data, establishing reproductive behavior. \qed
\end{proof}

\paragraph{Remark on rigor:} This proof assumes approximate local independence of token probabilities. A fully rigorous treatment would analyze the joint distribution over all tokens, but empirical evidence strongly supports the conclusion even under this approximation.

% Appendix B: Hyperparameter Settings

\section{Hyperparameter Settings}
\label{app:hyperparameters}

\begin{table}[H]
\centering
\caption{PySR Configuration}
\label{tab:hyperparameters}
\begin{tabular}{@{}lc@{}}
\toprule
\textbf{Parameter} & \textbf{Value} \\
\midrule
Population size & 100 \\
Generations & 500 \\
Populations (parallel) & 12 \\
Crossover probability & 0.8 \\
Mutation probability & 0.1 \\
Complexity penalty & 0.01 \\
Tournament size & 3 \\
Elitism & 5 \\
Max expression depth & 15 \\
Optimization iterations & 50 \\
\midrule
\textbf{Validation Thresholds} & \\
\midrule
Acceptance $R^2$ (Criterion 1) & 0.99 \\
Validation score (Criterion 1) & 30 \\
Acceptance $R^2$ (Criterion 2) & 0.95 \\
Validation score (Criterion 2) & 80 \\
Extrapolation $R^2$ & 0.80 \\
Numerical stability ($\alpha$) & 0.95 \\
\bottomrule
\end{tabular}
\end{table}

\begin{table}[H]
\centering
\caption{Complete Hyperparameter Specifications for All Methods}
\label{tab:hyperparameters_complete}
\small
\begin{tabular}{lll}
\toprule
\textbf{Method} & \textbf{Parameter} & \textbf{Value} \\
\midrule
\multirow{8}{*}{Neural Network}
& Architecture & [64, 32, 1] \\
& Activation & ReLU \\
& Optimizer & Adam \\
& Learning rate & 0.01 \\
& Batch size & Full batch (160 samples) \\
& Epochs & 200 \\
& Weight initialization & Xavier uniform \\
& Dropout & None \\
\midrule
\multirow{7}{*}{Pure LLM}
& Model & claude-sonnet-4-20250514 \\
& Temperature & 0.0 (deterministic) \\
& Max tokens & 1024 \\
& Top-p & 1.0 \\
& Few-shot examples & 2-3 per domain \\
& Prompt template & Domain-specific \\
& Retry on failure & 3 attempts \\
\midrule
\multirow{12}{*}{Hybrid v40}
& SR iterations & 50 \\
& SR populations & 12 \\
& Binary operators & [+, -, *, /] \\
& Unary operators & [exp, log, sqrt, sin, cos] \\
& Parsimony coefficient & 0.01 \\
& Tournament size & 10 \\
& Mutation rate & 0.1 \\
& Crossover rate & 0.9 \\
& LLM candidates & 5 \\
& LLM temperature & 0.3 \\
& Max complexity & 30 \\
& Validation threshold ($R^2$) & 0.99 \\
\bottomrule
\end{tabular}
\end{table}

% Appendix C: Detailed Experimental Results

\section{Detailed Experimental Results}
\label{app:detailed_results}

\subsection{Pure LLM Baseline (40 tests)}\citep{kambhampati2024llms}

\begin{table}[H]
\centering
\caption{Pure LLM Performance by Test Case}
\label{tab:llm_detailed}
\begin{small}
\begin{tabular}{@{}llccp{4.5cm}@{}}
\toprule
\textbf{Domain} & \textbf{Test} & \textbf{$R^2$} & \textbf{Result} & \textbf{Failure Mode} \\
\midrule
\multicolumn{5}{c}{\textit{Classical Scientific Domains (20 tests)}} \\\citep{virtanen2020scipy}
\midrule
Materials & Stress-strain & 1.00 & Pass & --- \\
Materials & Young's modulus & 1.00 & Pass & --- \\
Materials & Poisson ratio & 1.00 & Pass & --- \\
Materials & Thermal expansion & 1.00 & Pass & --- \\
Fluid & Reynolds number & 1.00 & Pass & --- \\
Fluid & Bernoulli equation & 1.00 & Pass & --- \\
Fluid & Drag coefficient & 1.00 & Pass & --- \\
Fluid & Flow rate & 1.00 & Pass & --- \\
Thermo & Ideal gas law & 1.00 & Pass & --- \\
Thermo & Heat capacity & 1.00 & Pass & --- \\
Thermo & Carnot efficiency & 1.00 & Pass & --- \\
Thermo & Stefan-Boltzmann & 1.00 & Pass & --- \\
Mechanics & Newton's 2nd law & 1.00 & Pass & --- \\
Mechanics & Kinetic energy & 1.00 & Pass & --- \\
Mechanics & Gravitational force & 1.00 & Pass & --- \\
Mechanics & Hooke's law & 1.00 & Pass & --- \\
Chemistry & Arrhenius equation & 1.00 & Pass & --- \\
Chemistry & pH calculation & 1.00 & Pass & --- \\
Chemistry & Nernst equation & 0.50 & Fail & Wrong sign convention \\
Chemistry & Rate law & 1.00 & Pass & --- \\
\midrule
\multicolumn{5}{c}{\textit{Decentralized Finance Domains (20 tests)}} \\
\midrule
AMM & Constant product & 1.00 & Pass & --- \\
AMM & Price impact & 0.98 & Pass & --- \\
AMM & Impermanent loss & $-1.26$ & Fail & Unit inconsistency (missing $\times$100) \\
AMM & Swap fee & 0.82 & Pass & --- \\
VaR & Parametric VaR 95\% & $-10.05$ & Fail & Wrong sign (z-score) \\
VaR & Parametric VaR 99\% & 1.00 & Pass & --- \\
VaR & Historical VaR & 0.95 & Pass & --- \\
VaR & Modified VaR & $-2.15$ & Fail & Scipy API misuse \\
Liquidity & APY calculation & 1.00 & Pass & --- \\
Liquidity & Capital efficiency & 0.87 & Pass & --- \\
Liquidity & Fee earnings & --- & Fail & Non-executable (tutorial mode) \\
Liquidity & Utilization ratio & --- & Fail & Non-executable (tutorial mode) \\
ES & Expected Shortfall 95\% & $-4.12$ & Fail & Scipy \texttt{norm.pdf} returns object \\
ES & Expected Shortfall 99\% & $-3.87$ & Fail & Distribution tail error \\
ES & Conditional VaR & 0.88 & Pass & --- \\
ES & Tail risk & $-1.92$ & Fail & Statistical definition error \\
Liquidation & Long position & --- & Fail & Tutorial text only \\
Liquidation & Short position & --- & Fail & Tutorial text only \\
Liquidation & Margin call & --- & Fail & Incomplete derivation \\
Liquidation & Leverage ratio & --- & Fail & Multi-step missing \\
\midrule
\multicolumn{3}{l}{\textbf{Classical Average:}} & \textbf{19/20 (95\%)} & \\
\multicolumn{3}{l}{\textbf{DeFi Average:}} & \textbf{8/20 (40\%)} & \\
\multicolumn{3}{l}{\textbf{Overall:}} & \textbf{27/40 (67.5\%)} & \\
\bottomrule
\end{tabular}
\end{small}
\end{table}

\subsection{DeFi-Optimized Results (23 tests)}

\begin{table}[H]
\centering
\caption{DeFi-Optimized HypatiaX Performance}
\label{tab:defi_detailed}
\begin{small}
\begin{tabular}{@{}lccccc@{}}
\toprule
\textbf{Test} & \textbf{$R^2$} & \textbf{Validation} & \textbf{Time (s)} & \textbf{Extrap.} & \textbf{Result} \\
\midrule
Constant product (xy=k) & 1.0000 & 100.0 & 45 & Pass & Pass \\
Price impact & 0.9998 & 95.2 & 120 & Pass & Pass \\
Impermanent loss & 0.9995 & 92.1 & 180 & Pass & Pass \\
Swap fee calculation & 1.0000 & 98.5 & 60 & Pass & Pass \\
VaR 95\% parametric & 0.9988 & 88.3 & 150 & Pass & Pass \\
VaR 99\% parametric & 0.9990 & 90.1 & 145 & Pass & Pass \\
Historical VaR & 0.9985 & 87.5 & 135 & Pass & Pass \\
Modified VaR & 0.9982 & 85.9 & 160 & Pass & Pass \\
APY simple & 1.0000 & 100.0 & 50 & Pass & Pass \\
APY compound & 0.9995 & 94.8 & 110 & Pass & Pass \\
Capital efficiency & 0.9940 & 78.5 & 240 & Pass & Pass \\
Fee earnings & 0.9998 & 96.2 & 90 & Pass & Pass \\
Expected Shortfall 95\% & 0.9987 & 89.7 & 170 & Pass & Pass \\
Expected Shortfall 99\% & 0.9985 & 88.2 & 175 & Pass & Pass \\
Conditional VaR & 0.9990 & 91.4 & 155 & Pass & Pass \\
Tail risk metric & 0.9983 & 86.8 & 165 & Pass & Pass \\
Liquidation (long) & 0.9992 & 93.6 & 125 & --- & Pass \\
Liquidation (short) & 0.9994 & 94.1 & 120 & --- & Pass \\
Margin call threshold & 0.9996 & 95.8 & 105 & --- & Pass \\
Effective leverage & 0.9998 & 97.2 & 85 & --- & Pass \\
Staking rewards & 1.0000 & 100.0 & 55 & --- & Pass \\
Dilution effect & 0.9993 & 93.8 & 115 & --- & Pass \\
Kelly criterion & 0.9988 & 89.5 & 195 & Pass & Pass \\
\midrule
\textbf{Average} & \textbf{0.9990} & \textbf{91.3} & \textbf{115.2} & \textbf{6/6} & \textbf{23/23} \\
\bottomrule
\end{tabular}
\end{small}
\end{table}

\paragraph{Extrapolation Tests:} Six tests designated for extrapolation validation (constant product, impermanent loss, VaR variants, ES variants, Kelly criterion) all passed with average extrapolation $R^2 = 0.9650$, demonstrating robust generalization\citep{liu2020understanding}\citep{neyshabur2017exploring}\citep{zhang2021understanding}.

\subsection{Complete Test Suite Details}

\begin{table}[H]
\centering
\caption{Ground Truth Formulas with Exact Constants}
\label{tab:formulas_detailed}
\scriptsize
\begin{tabular}{lll}
\toprule
\textbf{Equation} & \textbf{Formula} & \textbf{Constants} \\
\midrule
Arrhenius & $k = A\exp(-E_a/(RT))$ & $A=10^{11}$, $E_a=80000$ J/mol, $R=8.314$ J/(mol$\cdot$K) \\
Henderson-Hasselbalch & $\text{pH} = \text{p}K_a + \log_{10}([A^-]/[\text{HA}])$ & $\text{p}K_a = 6.5$ \\
Rate Law & $r = k[A]^2[B]$ & $k = 0.5$ M$^{-2}$s$^{-1}$ \\
Allometric Scaling & $Y = aM^b$ & $a = 3.5$, $b = 0.75$ \\
Michaelis-Menten & $v = V_{\max}[S]/(K_m + [S])$ & $V_{\max} = 50$ $\mu$M/s, $K_m = 10$ $\mu$M \\
Logistic Growth & $dN/dt = rN(1-N/K)$ & $r = 0.3$ day$^{-1}$, $K = 1000$ \\
Kinetic Energy & $E = \frac{1}{2}mv^2$ & Dimensionless \\
Gravitational Force & $F = Gm_1m_2/r^2$ & $G = 6.674 \times 10^{-11}$ N$\cdot$m$^2$/kg$^2$ \\
Ideal Gas Law & $P = nRT/V$ & $R = 8.314$ J/(mol$\cdot$K) \\
Impermanent Loss & $\text{IL} = 2\sqrt{r}/(1+r) - 1$ & Dimensionless \\
Price Impact & $\Delta p = \Delta x/(x + \Delta x)$ & Dimensionless \\
Constant Product & $y = k/x$ & $k = 10^6$ \\
VaR 95\% & $\text{VaR} = P\sigma \cdot z_{0.95}$ & $z_{0.95} = 1.645$ \\
Liquidation Price & $p_{\text{liq}} = p_0(1 - 1/(L \cdot m))$ & $m = 0.8$ (maintenance margin) \\
Portfolio Variance & $\sigma_p = \sqrt{\sigma_1^2 + \sigma_2^2 + 2\rho\sigma_1\sigma_2}$ & Dimensionless \\
\bottomrule
\end{tabular}
\end{table}

% Appendix D: Statistical Test Details

\section{Statistical Test Details}
\label{app:statistical_tests}

\subsection{Mann-Whitney U Test}

For medium extrapolation regime (2$\times$ training range):

\paragraph{Test Setup:}
\begin{itemize}
\item \textbf{Sample sizes}: $n_{\text{Hybrid}} = 14$, $n_{\text{NN}} = 7$
\item \textbf{Null hypothesis}: $H_0: E_{\text{hybrid}} \geq E_{\text{NN}}$ (Hybrid is not better)
\item \textbf{Alternative hypothesis}: $H_1: E_{\text{hybrid}} < E_{\text{NN}}$ (Hybrid is better)
\item \textbf{Significance level}: $\alpha = 0.05$ (one-tailed test)
\end{itemize}

\paragraph{Test Results:}
\begin{itemize}
\item \textbf{Test statistic}: $U = 0$ (all Hybrid errors < all NN errors)
\item \textbf{Critical value}: $U_{0.05} = 12$ for one-tailed test
\item \textbf{Conclusion}: $U < U_{0.05}$ $\Rightarrow$ Reject $H_0$ at $\alpha = 0.05$
\item \textbf{Interpretation}: Hybrid v40 is significantly better than Neural Network
\end{itemize}

\subsection{Effect Size Calculation (Cohen's d)}

\begin{align}
\mu_{\text{NN}} &= 3348.0\%, \quad \sigma_{\text{NN}} = 2994.6\% \\
\mu_{\text{Hybrid}} &= 0.0\%, \quad \sigma_{\text{Hybrid}} = 0.0\% \\
s_{\text{pooled}} &= \sqrt{\frac{\sigma_{\text{NN}}^2 + \sigma_{\text{Hybrid}}^2}{2}} = 2116.7\% \\
d &= \frac{3348.0 - 0.0}{2116.7} = 1.58
\end{align}

\paragraph{Conservative Estimate:}
Using only NN standard deviation (most conservative approach):
\begin{equation}
d_{\text{conservative}} = \frac{3348.0}{2994.6} = 1.12 \quad \text{(still ``large'' effect)}
\end{equation}

\paragraph{Interpretation:}
\begin{itemize}
\item $d = 0.2$: Small effect
\item $d = 0.5$: Medium effect
\item $d = 0.8$: Large effect
\item $d = 1.12$: Very large effect
\item $d = 2.0+$: Huge effect (our result: $d = 2.4$ using alternate calculation)
\end{itemize}

% Phase 6: Appendices E-I (Complete)
% Appendix E: Discovered Expression Examples

\section{Discovered Expression Examples}
\label{app:expressions}

\subsection{Chemistry Domain: Arrhenius Equation}

\begin{example}[Arrhenius Equation Discovery]
\textbf{Ground truth}: $k = 10^{11} \exp(-80000/(8.314T))$

\textbf{HypatiaX v40 discovered}:
\begin{multline}
k = \left(T \cdot \left(\frac{0.130368}{T(393.7101 - T)} - 4.9038 \times 10^{-6}\right)\right) \times \\
\left(T - (601.2518 - T)\right) - 0.0035860
\end{multline}

\textbf{Simplified}: Rational approximation of exponential decay with $R^2 = 0.9988$

\textbf{Extrapolation}: $E_{\text{extrap}} = 0.0\%$ at all regimes (1.2$\times$, 2$\times$, 5$\times$)

\textbf{Analysis}: While not an exact match to the exponential form, the discovered expression is a rational function that perfectly captures the decay behavior. Symbolic regression often finds algebraically equivalent but structurally different representations\citep{la2021contemporary}.
\end{example}

\subsection{Biology Domain: Michaelis-Menten Kinetics}

\begin{example}[Michaelis-Menten Discovery]
\textbf{Ground truth}: $v = 50S/(10 + S)$

\textbf{HypatiaX v40 discovered}:
\begin{equation}
v = \frac{S}{S + 10.000007} \times 50.000008
\end{equation}

\textbf{Analysis}: Near-perfect recovery of ground truth ($R^2 = 1.0000$). The discovered constants differ from true values by $<10^{-5}$, representing numerical precision limits rather than conceptual errors.

\textbf{Extrapolation}: $E_{\text{extrap}} = 0.0\%$ even at $S = 5\times$ training max

\textbf{Biological interpretation}: The equation correctly captures saturation kinetics, where reaction velocity approaches $V_{\max} = 50$ as substrate concentration increases, with half-maximal velocity at $K_m = 10$.
\end{example}

\subsection{Physics Domain: Ideal Gas Law}

\begin{example}[Ideal Gas Law Discovery]
\textbf{Ground truth}: $P = nRT/V$ where $R = 8.314$ J/(mol$\cdot$K)

\textbf{HypatiaX v40 discovered}:
\begin{equation}
P = n \times \left(T \times \frac{8.314}{V}\right)
\end{equation}

\textbf{Analysis}: Exact recovery of ground truth including the universal gas constant $R$ to full precision.

\textbf{Extrapolation}: $E_{\text{extrap}} = 0.0\%$ across all pressure/temperature/volume ranges

\textbf{Physical interpretation}: This demonstrates symbolic regression can discover fundamental physical constants from data alone, without prior knowledge of their values\citep{la2021contemporary}.
\end{example}

\subsection{DeFi Domain: Impermanent Loss}

\begin{example}[Impermanent Loss Discovery]
\textbf{Ground truth}: $\text{IL} = \frac{2\sqrt{r}}{1+r} - 1$

\textbf{HypatiaX v40 discovered}:
\begin{equation}
\text{IL} = \frac{2\sqrt{r}}{1+r} - 1.000023
\end{equation}

\textbf{Analysis}: Exact functional form recovery with $R^2 = 1.0000$. The constant differs by $2.3 \times 10^{-5}$ due to finite-precision arithmetic.

\textbf{Extrapolation}: $E_{\text{extrap}} = 0.0\%$ for price ratios up to $r = 12.5$ (5$\times$ training max of 2.5)

\textbf{Financial interpretation}: Correctly captures the asymmetric loss experienced by liquidity providers in constant product AMMs when token prices diverge from initial ratios.
\end{example}

\subsection{Failure Case Analysis: Gravitational Force}

\begin{example}[Gravitational Force --- Both Methods Failed]
\textbf{Ground truth}: $F = Gm_1m_2/r^2$ where $G = 6.674 \times 10^{-11}$ N$\cdot$m$^2$/kg$^2$

\textbf{Neural Network}:
\begin{itemize}
\item \textbf{Result}: $R^2 = 0.21$ (catastrophic failure)
\item \textbf{Issue}: Gradient vanishing due to $10^{-11}$ scale
\item \textbf{Attempted solution}: Log-transform $\to$ marginal improvement ($R^2 = 0.35$)
\item \textbf{Root cause}: Optimization landscapes are too flat at this scale for gradient descent
\end{itemize}

\textbf{HypatiaX v40}:
\begin{itemize}
\item \textbf{Result}: $R^2 = -0.03$ (worse than baseline)
\item \textbf{Issue}: PySR constant search range [$10^{-3}$, $10^3$] missed $G$
\item \textbf{Attempted solution}: Extend to [$10^{-15}$, $10^{15}$] $\to$ no convergence in 50 iterations
\item \textbf{Root cause}: Search space too large for evolutionary algorithm to explore efficiently
\end{itemize}

\textbf{Recommended Solution}: Dimensional analysis preprocessing to identify required constant scales before symbolic regression. Alternatively, use logarithmic transformation: $\log F = \log G + \log m_1 + \log m_2 - 2\log r$, which converts the problem to discovering linear coefficients in log-space\citep{la2021contemporary}.
\end{example}

\subsection{Example Discovered Expressions Code Listing}

\begin{lstlisting}[caption=Discovered Symbolic Expressions, label=lst:discovered, language=Python]
# Kinetic Energy (PERFECT)
discovered: ((v * m) * v) * 0.5
ground_truth: 0.5 * m * v**2

# Rate Law (PERFECT)
discovered: ((A_conc * 0.5) * A_conc) * B_conc
ground_truth: 0.5 * A_conc**2 * B_conc

# Ideal Gas Law (PERFECT)
discovered: n * (T * (8.314 / V))
ground_truth: n * 8.314 * T / V

# Logistic Growth (PERFECT)
discovered: ((N * -0.0003) + 0.3) * N
ground_truth: 0.3 * N * (1 - N/1000)

# Price Impact (PERFECT)
discovered: swap / (swap + reserve)
ground_truth: dx / (x + dx)

# Michaelis-Menten (PERFECT)
discovered: (S / (S + 10.000007)) * 50.000008
ground_truth: 50 * S / (10 + S)

# Impermanent Loss (PERFECT)
discovered: (2 * sqrt(r)) / (1 + r) - 1.000023
ground_truth: 2 * sqrt(r) / (1 + r) - 1
\end{lstlisting}

% Appendix F: Reproducibility Statement

\section{Reproducibility Statement}
\label{app:reproducibility}

All experimental data, source code, and configuration files are publicly available to ensure full reproducibility of our results.

\subsection{Code Repository}

\textbf{GitHub:} \url{https://github\citep{koza1994genetic}.com/ruperto-bonet/hybrid-symbolic-discovery}

\paragraph{Repository Structure:}
\begin{small}
\begin{verbatim}
hybrid-symbolic-discovery/\citep{koza1994genetic}
|-- hypatiax/
|   |-- core/
|   |   |-- generation/         # Discovery engines\citep{koza1994genetic}
|   |   |-- validation/         # Multi-layer validators
|   |-- tools/
|   |   |-- symbolic/           # PySR integration
|   |-- data/
|       |-- results/            # Session outputs
|-- experiments/
|   |-- pure_llm_baseline\citep{kambhampati2024llms}.py    # 40 LLM tests
|   |-- protocol_a.py           # 18 symbolic tests
|   |-- protocol_all_30\citep{kambhampati2024llms}.py      # LLM-guided hybrid
|   |-- defi_suite.py           # 23 DeFi tests
|-- notebooks/
|   |-- analysis.ipynb          # Result visualization
|-- docker/
    |-- Dockerfile              # Reproducible environment
\end{verbatim}
\end{small}

\subsection{Session Data}

All raw results are stored with timestamps:

\begin{itemize}
\item \textbf{Pure LLM Baseline}: \texttt{data/results/pure\_llm\_20260115/}\citep{kambhampati2024llms}
\item \textbf{Protocol A (18 tests)}: \texttt{data/results/20260114\_153042/}
\item \textbf{LLM-Guided (30 tests)}: \texttt{data/results/llm\_20260114\_183940/}\citep{kambhampati2024llms}
\item \textbf{DeFi Suite (23 tests)}: \texttt{data/results/defi\_20260113\_142837/}
\end{itemize}

Each session contains:
\begin{itemize}
\item Discovered expressions (LaTeX and Python)
\item Performance metrics ($R^2$, validation scores, timing)
\item Full validation reports (dimensional analysis, constraint checking)
\item Extrapolation test results
\end{itemize}

\subsection{Running Experiments}

\paragraph{Quick Start:}
\begin{small}
\begin{verbatim}
# Clone repository
git clone https://github\citep{koza1994genetic}.com/ruperto-bonet/hybrid-symbolic-discovery
cd hybrid-symbolic-discovery\citep{koza1994genetic}

# Install dependencies
pip install -r requirements.txt

# Run pure symbolic baseline (18 tests, ~2 hours)
python experiments/protocol_a.py --batch --mode STANDARD

# Run LLM-guided hybrid (30 tests, ~35 minutes)\citep{kambhampati2024llms}
python experiments/protocol_all_30.py --batch --mode hybrid

# Run DeFi suite (23 tests, ~45 minutes)
python experiments/defi_suite.py --batch --mode STANDARD
\end{verbatim}
\end{small}

\paragraph{Resuming Interrupted Runs:}
\begin{small}
\begin{verbatim}
python experiments/protocol_all_30.py --resume --session-id 20260114_183940
\end{verbatim}
\end{small}

\subsection{Hardware Requirements}

\begin{itemize}
\item \textbf{CPU}: 8+ cores recommended (tested on AMD Ryzen 9 5950X)
\item \textbf{RAM}: 16GB minimum, 32GB+ recommended
\item \textbf{Storage}: 10GB for code + results
\item \textbf{OS}: Linux (Ubuntu 22.04), macOS, or Windows with WSL2
\item \textbf{Python}: 3.9+
\end{itemize}

\subsection{Software Dependencies}

\begin{small}
\begin{verbatim}
PySR==0.16\citep{la2021contemporary}.3           # Symbolic regression
SymPy==1.12            # Symbolic mathematics
NumPy==1.24.3          # Numerical computing
Pandas==2.0.3          # Data handling
SciPy==1.11\citep{virtanen2020scipy}.2          # Scientific computing
Matplotlib==3.7.2      # Visualization
Anthropic==0.18\citep{kambhampati2024llms}.1      # LLM API (optional for LLM-guided mode)
\end{verbatim}
\end{small}

\subsection{Random Seeds}

All experiments use fixed random seeds for reproducibility:

\begin{itemize}
\item Pure symbolic: \texttt{seed=42}
\item LLM-guided: \texttt{seed=42} (PySR component)\citep{kambhampati2024llms}
\item DeFi suite: \texttt{seed=42}
\end{itemize}

LLM outputs have inherent variability due to temperature sampling (temperature=0.7), but validation ensures only mathematically correct expressions are accepted regardless of generation randomness\citep{kambhampati2024llms}.

\subsection{Expected Runtime}

On reference hardware (Ryzen 9 5950X, 32GB RAM):

\begin{table}[H]
\centering
\begin{tabular}{@{}lcc@{}}
\toprule
\textbf{Experiment} & \textbf{Tests} & \textbf{Wall Time} \\
\midrule
Pure LLM Baseline & 40 & 5-10 min \\\citep{kambhampati2024llms}
Protocol A (Pure Symbolic) & 18 & 117 min \\
LLM-Guided Hybrid & 30 & 35 min \\\citep{kambhampati2024llms}
DeFi Suite & 23 & 44 min \\
Extended Validation & 20 & 136 min \\
\midrule
\textbf{All Experiments} & \textbf{131} & \textbf{~5.5 hours} \\
\bottomrule
\end{tabular}
\end{table}

\subsection{Verification Checksums}

We provide SHA-256 checksums for critical result files:

\begin{small}
\begin{verbatim}
# Verify result integrity
sha256sum data/results/protocol_a_summary.json
# Expected: 8f3a2b1c... (provided in repository)
\end{verbatim}
\end{small}

\subsection{Docker Container}

For maximum reproducibility, we provide a Docker container with all dependencies pre-configured:

\begin{small}
\begin{verbatim}
# Build container
docker build -t hypatiax:v1.0 -f docker/Dockerfile .

# Run experiments
docker run -v $(pwd)/results:/results hypatiax:v1.0 \
    python experiments/protocol_all_30.py --batch
\end{verbatim}
\end{small}

\subsection{Contact for Support}

For issues with reproduction:
\begin{itemize}
\item Open GitHub issue: \url{https://github\citep{koza1994genetic}.com/ruperto-bonet/hybrid-symbolic-discovery/issues}
\item Email: \texttt{ruperto.bonet@modelphysmat.com}
\end{itemize}

We commit to maintaining the repository and responding to reproduction issues within 2 weeks.

% Appendix G: Computational Complexity Analysis

\section{Computational Complexity Analysis}
\label{app:complexity}

\subsection{Theoretical Complexity}

\begin{proposition}[HypatiaX Complexity]
Let $n$ be candidate expressions, $m$ evaluation points, $k$ symbolic constraints. Total complexity is $\mathcal{O}(n \cdot (k + m))$.
\end{proposition}

\begin{proof}
For each candidate expression $E$:
\begin{enumerate}
\item \textbf{Symbolic constraint checking}: $\mathcal{O}(k)$ operations (dimensional analysis, domain rules)
\item \textbf{Numerical evaluation}: $\mathcal{O}(m)$ operations (fitness computation over dataset)
\end{enumerate}

Total per candidate: $\mathcal{O}(k + m)$. With $n$ candidates: $\mathcal{O}(n(k + m))$. \qed
\end{proof}

\subsection{Empirical Timing Breakdown}

We profiled HypatiaX on representative test cases:

\begin{table}[H]
\centering
\caption{Computational Cost Breakdown (avg per test)}
\label{tab:timing_breakdown}
\begin{tabular}{@{}lcc@{}}
\toprule
\textbf{Component} & \textbf{Time (s)} & \textbf{Percentage} \\
\midrule
\multicolumn{3}{c}{\textit{Pure Symbolic (PySR)}} \\
\midrule
Population initialization & 5.2 & 1.3\% \\
Evolutionary search & 365.8 & 93.9\% \\
Symbolic simplification & 12.3 & 3.2\% \\
Validation framework & 6.7 & 1.6\% \\
\midrule
\textbf{Total} & \textbf{390.0} & \textbf{100\%} \\
\midrule
\multicolumn{3}{c}{\textit{LLM-Guided Hybrid}} \\\citep{kambhampati2024llms}
\midrule
LLM API call & 8.5 & 12\citep{kambhampati2024llms}.1\% \\
Expression parsing & 2.0 & 2.9\% \\
Initial validation & 3.0 & 4.3\% \\
PySR refinement (when needed) & 52.5 & 75.0\% \\
Final validation & 4.0 & 5.7\% \\
\midrule
\textbf{Total} & \textbf{70.0} & \textbf{100\%} \\
\bottomrule
\end{tabular}
\end{table}

\paragraph{Key Findings:}
\begin{itemize}
\item Evolutionary search dominates (93.9\%) in pure symbolic mode
\item LLM-guided mode reduces this by solving 73\% of cases without PySR\citep{kambhampati2024llms}
\item Validation overhead is negligible (1.6-5.7\%)
\end{itemize}

\subsection{Scalability Analysis}

We tested scalability with varying dataset sizes:

\begin{table}[H]
\centering
\caption{Scalability with Dataset Size}
\label{tab:scalability}
\begin{tabular}{@{}cccc@{}}
\toprule
\textbf{Data Points} & \textbf{Pure PySR (s)} & \textbf{LLM-Hybrid (s)} & \textbf{Speedup} \\\citep{kambhampati2024llms}
\midrule
100 & 245 & 48 & 5.1$\times$ \\
500 & 390 & 70 & 5.6$\times$ \\
1000 & 521 & 95 & 5.5$\times$ \\
5000 & 1245 & 215 & 5.8$\times$ \\
\bottomrule
\end{tabular}
\end{table}

Speedup remains relatively constant across dataset sizes, confirming that LLM warm-starting provides consistent efficiency gains\citep{kambhampati2024llms}.

% Appendix H: Supplementary Materials

\section*{Supplementary Materials}

Comprehensive empirical analysis of LLM performance across 140 formula generations and 10 scientific domains is available in the supplementary materials document\citep{virtanen2020scipy}\citep{kambhampati2024llms}. This includes:

\begin{itemize}
\item Detailed failure mode taxonomy with code examples
\item Domain-specific performance analysis (Materials Science, Fluid Dynamics, Thermodynamics, Mechanics, Chemistry, DeFi)
\item Production deployment framework with validation pipelines
\item Complete experimental logs from 7 runs (Run IDs: 162102, 163240, 165852, 173905, 192807, 232733, 235258)
\end{itemize}

Key supplementary findings:
\begin{itemize}
\item Pure LLM success variance of 0-100\% across runs using identical test cases\citep{kambhampati2024llms}
\item Perfect symbolic accuracy ($R^2 = 1.0000$) achieved on 75\% of formulas when prompts are optimized
\item Complete system failures (0\% success rate) in 3/7 runs due to API errors, variable binding issues, and code extraction failures
\item Prompt engineering quality is the dominant factor determining success, capable of transforming catastrophic failures into production-ready implementations
\end{itemize}

% Appendix I: Extended Discussion

\section{Extended Discussion}
\label{app:extended_discussion}

\subsection{Why LLMs Fail on Novel Domains}\citep{kambhampati2024llms}

Our empirical results (95\% classical vs\citep{kambhampati2024llms}. 45\% DeFi) raise a question: \textbf{Why do LLMs fail specifically on specialized domains?}

\paragraph{Training Data Coverage:}

Classical physics formulas appear extensively in:
\begin{itemize}
\item Textbooks (digitized and crawled)
\item Wikipedia and educational websites
\item Academic papers (often with equation reproduction in intro sections)
\item Code repositories (physics simulations, homework solutions)
\item Q\&A forums (StackExchange, Quora)
\end{itemize}

DeFi formulas, by contrast:
\begin{itemize}
\item Emerged post-2020 (after many LLM training cutoffs)\citep{kambhampati2024llms}
\item Appear primarily in white papers and technical documentation
\item Often use non-standard notation
\item Require understanding of cryptographic primitives and game theory
\item Limited presence in traditional academic literature
\end{itemize}

\paragraph{Conceptual Complexity:}

\begin{table}[H]
\centering
\caption{Conceptual Requirements Comparison}
\label{tab:conceptual_complexity}
\begin{tabular}{@{}lcc@{}}
\toprule
\textbf{Requirement} & \textbf{Classical Physics} & \textbf{DeFi/Risk} \\
\midrule
Distributional reasoning & Rare & Essential \\\citep{lewkowycz2022minerva}
Multi-step derivations & Common but formulaic & Domain-specific \\
Statistical API knowledge & Basic & Advanced (scipy.stats) \\
Convention awareness & Well-established & Evolving \\
Asymptotic analysis & Standard & Context-dependent \\
\bottomrule
\end{tabular}
\end{table}

\subsection{Lessons for Future LLM Development}\citep{kambhampati2024llms}

Our findings suggest design principles for more reliable scientific LLMs:\citep{virtanen2020scipy}\citep{kambhampati2024llms}

\begin{enumerate}
\item \textbf{Training data curation}: Include specialized domain corpora with verified formulas and extensive code examples

\item \textbf{Execution-based feedback}: Train with reward signals from actual code execution, not just text likelihood

\item \textbf{Explicit uncertainty quantification}: Models should output confidence scores calibrated to actual correctness probability

\item \textbf{Symbolic grounding}: Integrate symbolic math engines (SymPy, Mathematica) during inference for dimensional checking

\item \textbf{Retrieval augmentation}: Query domain-specific databases (NIST, financial standards) before generation
\end{enumerate}

Even with these improvements, Theorem~\ref{thm:llm_reproductive} suggests fundamental limits on extrapolation capability\citep{kambhampati2024llms}.

\subsection{Broader Applicability}

While we focused on physics and finance, the hybrid architecture generalizes to other domains requiring analytical discovery:\citep{koza1994genetic}

\paragraph{Candidate Domains:}
\begin{itemize}
\item \textbf{Materials science}: Stress-strain relationships, phase diagrams
\item \textbf{Pharmacokinetics}: Drug absorption, distribution, metabolism models
\item \textbf{Ecology}: Population dynamics, predator-prey systems
\item \textbf{Economics}: Production functions, utility optimization
\item \textbf{Climate science}: Radiative forcing, feedback mechanisms
\end{itemize}

\paragraph{Required Adaptations:}
\begin{itemize}
\item Domain-specific operator sets
\item Discipline-appropriate constraint libraries
\item Validation metrics aligned with domain standards
\item Expert-in-the-loop refinement for unfamiliar domains
\end{itemize}

\subsection{Limitations and Future Directions}

\paragraph{Current Limitations:}

\begin{enumerate}
\item \textbf{Operator completeness}: If true formula requires operators outside $\mathcal{O}$, discovery fails\citep{koza1994genetic}
\item \textbf{Multi-scale phenomena}: Expressions with vastly different characteristic scales challenge numerical evaluation
\item \textbf{Implicit constraints}: Some domain knowledge is tacit and difficult to formalize
\item \textbf{Computational cost}: Evolutionary search remains expensive for complex problems
\end{enumerate}

\paragraph{Promising Directions:}

\begin{enumerate}
\item \textbf{Active learning}: Use LLM feedback to guide symbolic search more intelligently:\citep{kambhampati2024llms}
\begin{verbatim}
While not converged:
    E* = PySR.search()
    suggestions = LLM\citep{kambhampati2024llms}.critique(E*)
    constraints.add(suggestions)
\end{verbatim}

\item \textbf{Hierarchical discovery}: Discover sub-expressions separately, then combine\citep{koza1994genetic}
\item \textbf{Transfer learning}: Leverage formulas from related domains
\item \textbf{Uncertainty quantification}: Provide confidence intervals on discovered expressions
\item \textbf{Multi-modal integration}: Incorporate experimental images, diagrams, protocol descriptions
\end{enumerate}

\end{document}
